
\section{The non-star part of the character expansion and coincident labels}\label{sec:extraction}

\paragraph{What this section proves and how to read it.}
Expanding a cut kernel in characters gives terms indexed by sets of
cross pairs. Terms whose pairs share an endpoint (\emph{stars}) are
handled in \cref{sec:spectral}; this section bounds the remaining
terms and the terms where labels coincide, showing both are second
order with explicit constants. No new cancellation is proved here:
the work is bookkeeping to keep constants small when many subsets
are summed. \emph{Extraction} means the following choice, made for
each group of subsets (a \emph{class} $\mathcal C(P,A)$: pairs in
$P$ present, pairs in $A$ absent, others free): choose two disjoint
present pairs (\emph{witnesses}, certifying second order), choose
sets of vertices meeting every other undecided pair (a \emph{vertex
cover}, whose labels are conditioned on), and choose labelling
orders. An \emph{extraction datum} records these choices (called a
``record of a class'' in the text below, not to be confused with
records of Part~III). Inputs: \cref{sec:cuts,sec:dl}. Output: the
per-class bound (\cref{subsec:extraction:extraction}) used by the
spectral criteria. A first reading should follow the class
definition and the statement of the per-class bound; the counting
details are reference material. Writing each cross constraint at $e=(j,k)$ as
$\tfrac12(1+\chi(\varsigma_{jk}(x'_k-x_j)))$ gives the expansion; the non-star part
$T_{\rm ns}$ is second order because a non-star contains two
disjoint pairs, and the coincidence part $T_{\rm coll}$ corrects for
half labellings sharing a label. The difficulty is the constant at
the first prime of a range, where $2^{m_\times}$ subsets are summed;
it is kept small by stopping at the first two disjoint pairs,
conditioning only on a vertex cover, counting constraints met by
each vertex, and labelling the coincident vertices (those forced
to carry equal labels) first in the coincidence count.

\subsection{The character expansion and the compressed partition}\label{subsec:extraction:partition}

\paragraph{Setting.}
Throughout this section, except in \cref{subsec:push:pinnedremainder},
$(G,\vartheta,U^+,\mathfrak o)$ are spectral data (\cref{def:spec:data}) with the
notation fixed there: $n=|U^+|$, the half graph $G^+$, adjacency $j\sim k$ in
$G^+$, the degrees $d^0_k$ and $\Delta^+=\max_kd^0_k$, the cross signs
$\varsigma_{jk}$, the neighbourhoods $N(k)=\{j:\varsigma_{jk}\ne0\}$ with
$r_k=|N(k)|$, $d_k=d^0_k+r_k$, $\Delta=\max_kd_k$, the cross pairs $E_\times$ and
$m_\times=|E_\times|$. We write $\operatorname{Adj}(i)=\{j\in U^+:j\sim i\}$, so
that $|\operatorname{Adj}(i)|=d^0_i$, and $e(L)$, $e(L,L')$ for the numbers of
edges of $G^+$ inside $L$ and between disjoint sets $L,L'$. The prime $p$ is
odd, $r=p^{-1/2}$, $e_t(p,n')=2^t(h_tr+n'r^2)$ is the count error of
\eqref{eq:graphs:et}, $\mathcal I^+=\operatorname{Inj}_p(G^+)$ and
$\mu=\lambda_p(G^+)$ is the half law. For $Y\subseteq U^+$ put
\[
 R_Y=\prod_{i\in Y}\frac{1+e_{d^0_i}(p,n)}{1-e_{d^0_i}(p,n)},\qquad
 R_{\rm full}=R_{U^+},\qquad \hat e=e_{\Delta^+}(p,n),
\]
the completion ratios of \cref{lem:graphs:vertex} for the graph $G^+$. Under
(S1) at $p$, $e_t(p,n')<1$ for $t\le\Delta$ and $n'\le2n$; in particular
$\hat e<1$, and every factor of every $R_Y$ is defined and at least one.

The two labellings of a pair $(x,x')\in\mathcal I^+\times\mathcal I^+$ live on two
copies $X$ and $Z$ of $U^+$; in a cross pair $(j,k)$ the $x$-endpoint $j$ is a
vertex of $X$ and the $x'$-endpoint $k$ a vertex of $Z$. Graph quantities
($\operatorname{Adj}$, $d^0$, $N$, $r$) of a vertex of $X$ or $Z$ are those of
the corresponding vertex of $U^+$. The characters
$\chi_e(x,x')=\chi(\varsigma_{jk}(x'_k-x_j))$ of the cross pairs $e=(j,k)$, their
products $\chi_S$, the operators $T_S$ on $L^2(\mu)$ with kernel $\chi_S$ relative to
$\mu\otimes\mu$, the stars and nonstars, $T_{\rm ns}=\sum_{S\ \rm nonstar}T_S$, and
the collision operator $T_{\rm coll}$, whose kernel $K_{\rm coll}$ is
$\prod_{e\in E_\times}(1+\chi_e)$ on the pairs $(x,x')$ that share a label and $0$
elsewhere, are those of \cref{subsec:spec:perp}. By \cref{lem:spec:expansion},
$2^{m_\times}$ times the cut kernel is
$\one\otimes\one+\sum_{S\ \rm star}T_S+T_{\rm ns}-T_{\rm coll}$
\eqref{eq:record:root-kernel}, and $K_{\rm coll}$ is nonnegative and symmetric; the
constants and the stars are treated in \cref{sec:spectral}. By \cref{lem:spec:elementary},
$(j,k)\in E_\times$ iff $(k,j)\in E_\times$; so $j\in N(k)$ iff $k\in N(j)$, and $r_j$
is also the number of cross pairs with $x$-endpoint $j$.

\paragraph{Graded errors.}
Every statement of this section outside \cref{subsec:push:pinnedremainder}, whose
rows carry their own error $e^{\rm pin}_d$, holds with the same proof when the
coarse error $e_t(p,n')$ is replaced throughout (in $R_Y$, $R_{\rm full}$, $\hat e$,
the vertex and leaf factors, $L_{\rm int}$, $\Upsilon$ and $\Upsilon_k$) by the graded
error $e^-_t(n')$ of \eqref{eq:graphs:graded}. Indeed the proofs use three
properties of $e_t(p,n')$: it bounds the relative error of a count of
\cref{lem:extension} from above whatever the excluded set, and from below when at
most $n'$ labels are excluded
(\cref{lem:extension}\ref{it:ext:upper},\ref{it:ext:coarse}); it is nondecreasing in
$t$ and $n'$ and a polynomial in $r$ with nonnegative coefficients
(\ref{it:et:poly}); and the densities and completion ratios of
\cref{lem:graphs:vertex} are formed with it. The graded error has the first two
properties by \cref{lem:extension}\ref{it:ext:upper},\ref{it:ext:lower} with
\ref{it:et:compare} and \ref{it:et:poly}. For the third,
\cref{lem:graphs:vertex}\ref{it:vertex:dens}--\ref{it:vertex:tuple} hold with
$e^-_t(n)$ in place of $e_t(p,n)$, under $e^-_\Delta(n)<1$: in their proofs
\cref{lem:extension}\ref{it:ext:coarse} enters only through lower counts, for which
\ref{it:ext:lower} with $z=n$ serves, and
\cref{lem:graphs:sequential}\ref{it:seq:count},\ref{it:seq:coarse} for the coarse pair
may be replaced by \ref{it:seq:count} for the graded pair, whose ratio
$\prod_i(1+e^+_{t_i})/(1-e^-_{t_i}(n))$ is at most
$\prod_i(1+e^-_{d_i}(n))/(1-e^-_{d_i}(n))$ by \ref{it:et:compare} and
\ref{it:et:poly}, since $t_i\le d_i$. Where a proof uses the explicit form of
$e_t(p,n')$ (\cref{cor:extraction-compressed}), it applies through
$e^-_t(n')\le e_t(p,n')$ (\ref{it:et:compare}); in particular, under (S1),
$e^-_t(n')<1$ for $t\le\Delta$ and $n'\le2n$. The \emph{error pair of the setting}
is the coarse pair or the graded pair of \cref{def:graphs:error-pair}; one choice
is made for all the errors of a bound, and the nonstar coefficients of the root cuts
of the upper range below $2^{52}$ use the graded one.

\paragraph{Classes.}
For disjoint $P,A\subseteq E_\times$ the \emph{class}
$\mathcal C(P,A)=\{S:P\subseteq S\subseteq E_\times\setminus A\}$ consists of the
sets in which the pairs of $P$ are \emph{present}, those of $A$ \emph{absent}, and
those of $F=E_\times\setminus(P\cup A)$ \emph{free}. A pair is \emph{decided} if it
is present or absent; pairs of $X\times Z$ outside $E_\times$ count as absent. The
\emph{mass} of a class is $2^{|F|}$, and $T_{\mathcal C}=\sum_{S\in\mathcal C}T_S$.
An \emph{exact partition} of a family of subsets of $E_\times$ is a set of
pairwise disjoint classes whose union is the family. If $\mathcal C$ is the
disjoint union of $\mathcal C_1$ and $\mathcal C_2$, then
$T_{\mathcal C}=T_{\mathcal C_1}+T_{\mathcal C_2}$. Hence bounds for the classes
of an exact partition add up to a bound for the whole family, and for a class
split in two the smaller of a bound for the class and the sum of bounds for its
two parts is again a bound for the class.

\begin{definition}[Witnesses and covers]\label{def:extraction:witness}
Let $\mathcal C=\mathcal C(P,A)$ be a class with free pairs $F$, and $k\ge1$.
\begin{enumerate}[label=\textup{(\alph*)}]
\item A \emph{witness of size $k$} consists of pairs
 $(j_1,k_1),\dots,(j_k,k_k)\in P$ with distinct $j_i$ and distinct $k_i$ such
 that, with $V=\{j_1,\dots,j_k\}\subseteq X$ and $W=\{k_1,\dots,k_k\}\subseteq Z$,
 every pair of the \emph{rectangle} $E_{VW}=(V\times W)\cap E_\times$ is decided;
 call this \textup{(W$_k$)}.
\item The \emph{witness digraph} $D_{\mathcal C}$ has the vertex set
 $\{1,\dots,k\}$ and an arc $i\to l$ ($i\ne l$) whenever $(j_i,k_l)\in P$. If
 all its strongly connected components have at most two vertices,
 $\beta(\mathcal C)$ denotes the number of components with two vertices.
\item A \emph{cover} is a pair $C_x\subseteq X\setminus V$, $C_z\subseteq
 Z\setminus W$ such that every pair of $P\cup F$ outside $V\times W$ has its
 $x$-endpoint in $C_x$ or its $x'$-endpoint in $C_z$; call this
 \textup{(VC$_k$)}. The \emph{integrated} vertices are
 $U_x=X\setminus(V\cup C_x)$ and $U_z=Z\setminus(W\cup C_z)$. The pair
 $(X\setminus V,Z\setminus W)$ is always a cover.
\end{enumerate}
\end{definition}

The partition lemma concerns only the combinatorics of the cross pairs, and we
state it for an arbitrary finite set $E_\times\subseteq X\times Z$ of pairs,
symmetric or not (the pinned rows of \cref{subsec:push:pinnedremainder} need a
set that is not symmetric). For such a set let $r^X_j$ and $r^Z_k$ be the numbers
of its pairs with $x$-endpoint $j$ and with $x'$-endpoint $k$; for spectral data
both are the cross degrees $r_j$, $r_k$. Two pairs are \emph{disjoint} if their
$x$-endpoints differ and their $x'$-endpoints differ.

\begin{lemma}[The compressed partition]\label{lem:advance:compressed-partition}
Let $E_\times\subseteq X\times Z$ be finite, $m_\times=|E_\times|$, with a total
order $\prec$.
\begin{enumerate}[label=\textup{(\alph*)}]
\item A subset of $E_\times$ contains no two disjoint pairs if and only if it is
 empty or a star. The empty set and the stars are
 \[
  1+m_\times+\sum_j\bigl(2^{r^X_j}-1-r^X_j\bigr)+\sum_k\bigl(2^{r^Z_k}-1-r^Z_k\bigr)
 \]
 subsets; so the masses of the classes of an exact partition of the nonstar
 subsets add up to $2^{m_\times}$ minus this number.
\item For a nonstar $S$ let $e=(a,b)$ be its first pair. Exactly one of the
 following holds, and it defines a class $\mathcal C(S)=\mathcal C(P,A)$.
 \begin{enumerate}[label=\textup{(\Alph*)}]
 \item Some pair of $S$ is disjoint from $e$. Let $f=(a',b')$ be the first one,
  let $\mathrm{Con}=\{(a,b'),(a',b)\}\cap E_\times$ be the \emph{connectors}, and
  \[
   P=\{e,f\}\cup(S\cap\mathrm{Con}),\qquad
   A=\{g:g\prec e\}\cup\{g\prec f:g\text{ disjoint from }e\}\cup(\mathrm{Con}\setminus S).
  \]
 \item No pair of $S$ is disjoint from $e$. Then $S$ contains, besides $e$, a pair
  with $x$-endpoint $a$ and a pair with $x'$-endpoint $b$ (the \emph{arms} of $e$).
  Let $g_1=(a,b'')$ and $g_2=(a'',b)$ be the first arms of the two kinds, let $A_1$
  ($A_2$) be the set of pairs $g\ne e$ with $x$-endpoint $a$ ($x'$-endpoint $b$) and
  $g\prec g_1$ ($g\prec g_2$), and
  \[
   P=\{e,g_1,g_2\},\qquad
   A=\{g:g\prec e\}\cup A_1\cup A_2\cup\{g:g\text{ disjoint from }e\}.
  \]
 \end{enumerate}
 Every element of $\mathcal C(S)$ is a nonstar whose class is $\mathcal C(S)$.
 So the distinct classes $\mathcal C(S)$ form an exact partition of the nonstar
 subsets, the \emph{compressed partition}. In case \textup{(A)} the pairs $e,f$,
 and in case \textup{(B)} the arms $g_1,g_2$, form a witness of size two of
 $\mathcal C(S)$. Its witness digraph is a two-cycle if, in case \textup{(A)},
 both connectors are present, and it is acyclic otherwise.
\item Let $\mathcal C(P,A)$ be a class and $g$ a free pair. The \emph{children}
 $\mathcal C(P\cup\{g\},A)$ and $\mathcal C(P,A\cup\{g\})$ partition it; so
 replacing a class of an exact partition by its children, any number of times,
 gives an exact partition. A witness of $\mathcal C(P,A)$ is a witness of each
 child, with the same witness digraph, and a cover of $\mathcal C(P,A)$ is a
 cover of each child. If $k$ pairs of $P$ have distinct $x$-endpoints and
 distinct $x'$-endpoints, they form a witness of size $k$ of every class
 obtained from $\mathcal C(P,A)$ by deciding, one at a time, the undecided pairs
 of their rectangle.
\end{enumerate}
\end{lemma}

\begin{proof}
(a) Let $S$ be nonempty without two disjoint pairs, and fix $e=(a,b)\in S$.
Every other pair of $S$ shares $a$ or $b$ with $e$. If $S$ contained a pair
$(a,b_1)$ with $b_1\ne b$ and a pair $(a_1,b)$ with $a_1\ne a$, these two would be
disjoint. So all pairs of $S$ have the $x$-endpoint $a$, or all have the
$x'$-endpoint $b$. Conversely any two pairs of a star share an endpoint. A set of
at least two pairs is not both a row star and a column star, since two pairs with
the same two endpoints coincide, while a single pair is both. Counting the empty
set, the $m_\times$ single pairs, the row stars of at least two pairs inside each row
and the column stars of at least two pairs inside each column gives the number.

(b) By (a) a nonstar $S$ contains two disjoint pairs, and so it has a first pair
$e$ and falls in exactly one of the two cases. In case (B) every pair of $S$
shares an endpoint with $e$, and $S$ is neither contained in the row of $e$ nor
in its column, which gives the arms; $g_1$ and $g_2$ are disjoint, and both
follow $e$. In both cases $P\cap A=\emptyset$, $P\subseteq S$, and no pair of $A$
lies in $S$: the pairs before $e$ by the choice of $e$, the pairs disjoint from
$e$ before $f$ by the choice of $f$, the arms before $g_1$, $g_2$ by the choice of
the arms, and in case (B) the pairs disjoint from $e$ by assumption. So
$S\in\mathcal C(S)$. Let $S'\in\mathcal C(S)$. In case (A), $e\in S'$ and no pair
before $e$ lies in $S'$, so $e$ is the first pair of $S'$; $f\in S'$ is disjoint
from $e$ and no pair disjoint from $e$ precedes it in $S'$; and
$S'\cap\mathrm{Con}=S\cap\mathrm{Con}$. Hence $S'$ is a nonstar in case (A) with the
same $e$, $f$ and connectors, and $\mathcal C(S')=\mathcal C(S)$. In case (B), $e$
is again the first pair of $S'$, no pair of $S'$ is disjoint from $e$, and $g_1$,
$g_2$ are the first arms of $S'$; $S'$ is a nonstar because it contains the
disjoint pairs $g_1,g_2$, and $\mathcal C(S')=\mathcal C(S)$. So two classes
$\mathcal C(S_1)$, $\mathcal C(S_2)$ with a common element coincide, and the
classes cover the nonstar subsets exactly once.

In case (A) take $(j_1,k_1)=e$ and $(j_2,k_2)=f$, so $V=\{a,a'\}$ and
$W=\{b,b'\}$. The rectangle consists of $e$, $f$ and the connectors, all decided;
$D_{\mathcal C}$ has the arc $1\to2$ iff $(a,b')\in P$ and the arc $2\to1$ iff
$(a',b)\in P$. In case (B) take $(j_1,k_1)=g_1=(a,b'')$ and $(j_2,k_2)=g_2=(a'',b)$,
so $V=\{a,a''\}$ and $W=\{b'',b\}$. The rectangle consists of $g_1$, $g_2$, the
present pair $e=(a,b)$, which gives the arc $1\to2$, and the pair $(a'',b'')$,
which is disjoint from $e$ and therefore absent; there is no arc $2\to1$.

(c) The children are the elements of $\mathcal C(P,A)$ that contain $g$ and those
that do not. A child has more present and more absent pairs; the witness pairs
stay present, the rectangle stays decided, and $P\cap(V\times W)$ does not change,
because $g$ was free and so lies outside the rectangle. The non-absent pairs of a
child are non-absent in the parent, so (VC$_k$) passes to the child. Deciding the
undecided pairs of a rectangle produces classes in which the rectangle is
decided, which is (W$_k$).
\end{proof}

\begin{remark}[Comments on the partition]\label{rem:extraction-partition}
(1) The arm case cannot be omitted: a three-edge path whose middle pair comes
first contains no pair disjoint from that pair. In the arm case the first pair $e$
is a present connector of the witness, not a witness pair.
(2) A set $J\times K$ with $|J|,|K|\ge2$ contains two disjoint pairs, so complete
bipartite subsets lie in classes and need no separate estimate.
(3) Refinement does not worsen the available bound: by (c) the witness and cover of a class remain valid at
all its descendants, so the size-two bound is available for every class obtained by
repeated refinement, and such a class may keep the smaller of its own bound
and the sum of the bounds of its children.
Larger witnesses are found along augmenting paths: if the present pairs contain a
matching of size $k+1$, the symmetric difference with a matching $M$ of size $k$
contains an augmenting path for $M$ (Berge's lemma), of length one, three or five
when $k=2$. The search for witnesses is heuristic and lies outside the proof; a
proposed witness is used only after \textup{(W$_k$)} and the digraph
condition have been verified.
(4) Summed over a refinement of the compressed partition, the bounds for the
classes at which refinement stops (the leaves) give the
nonstar coefficients $C_k(Q)$ of \cref{cor:extraction-leaf}(b). Every choice (the
order, the pair to decide, the witness, when to stop, the cover and the orders)
affects only constants. The certificates make the choices once, in
floating point at one prime, and fix them for all primes; their builders also check
that the masses of the classes and the number in (a) add up to $2^{m_\times}$, and
that a refined union of classes carries the mass it replaces. These checks accompany
the proof; they do not replace it. (In the data files, whose
names are reproduced in \cref{app:certificates}, the classes of
case~(B) are called type~A and those of case~(A) type~B --- the
reverse of the letters in this text; the cases here are always
referred to as (A) and (B).)
\end{remark}

\subsection{Bounding a class by two designated pairs}\label{subsec:extraction:kernel}

Fix a class $\mathcal C(P,A)$ and its designated witness pairs
(\cref{def:extraction:witness}). After conditioning on all other
labels, the kernel that remains on the witness labels is a product
of characters $\chi(\pm(y_l-x_i))$ over a set $\mathcal K$ of
pairs. The lemma bounds such a kernel: its compressed norm is at most $2^\beta r^k$
($k$ the number of witness labels, so second order as soon as
$k\ge2$), provided the directed graph of the pairs has no strongly
connected component larger than a $2\times2$ rectangle, with $\beta$ the
number of two-vertex components. These components use the four-cycle
bound of \cref{sec:dl}; larger components require a separate estimate.

\begin{lemma}[Witness kernels]\label{lem:matching:kernel}
Let $k\ge1$, let $\mathcal K\subseteq\{1,\dots,k\}^2$ contain the diagonal pairs
$(i,i)$, let $\varsigma_{il}\in\{\pm1\}$ for $(i,l)\in\mathcal K$, and let
$T_{\mathcal K}$ be the operator on $L^2(\Fp^k)$, with uniform probability, given
by
\[
 (T_{\mathcal K}f)(x)=\E_{y\in\Fp^k}\Bigl[\prod_{(i,l)\in\mathcal K}
 \chi\bigl(\varsigma_{il}(y_l-x_i)\bigr)\,f(y)\Bigr].
\]
Let $D_{\mathcal K}$ be the digraph on $\{1,\dots,k\}$ with an arc $i\to l$ for
every $(i,l)\in\mathcal K$ with $i\ne l$.
\begin{enumerate}[label=\textup{(\alph*)}]
\item If $D_{\mathcal K}$ is acyclic, then $\norm{T_{\mathcal K}}\le r^k$.
\item Suppose that every strongly connected component of $D_{\mathcal K}$ has at
 most two vertices and that $\beta$ of them have two. Let $P_{\rm in}$ and
 $P_{\rm out}$ be the multiplications by the indicators of sets of tuples $y$,
 respectively $x$, with $y_i\ne y_l$, respectively $x_i\ne x_l$, for every
 two-vertex component $\{i,l\}$ (for instance the injective tuples). Then
 $\norm{P_{\rm out}T_{\mathcal K}P_{\rm in}}\le2^\beta r^k$.
\item $D_{\mathcal K}$ is acyclic if and only if the identity is the only
 permutation $\pi$ of $\{1,\dots,k\}$ with $(i,\pi(i))\in\mathcal K$ for
 all $i$, that is, if and only if the diagonal is the only perfect matching
 contained in $\mathcal K$.
\item \emph{Inserted multipliers.} Let $k=2$ and $\mathcal K=\{(1,1),(2,2)\}$, and
 let $F_{\rm out}(x_1,x_2)$, $F_{\rm mid}(x_1,y_2)$ and $F_{\rm in}(y_1,y_2)$ be
 bounded functions. The operator on $L^2(\Fp^2)$ with kernel
 \[
  \chi\bigl(\varsigma_{11}(y_1-x_1)\bigr)\,\chi\bigl(\varsigma_{22}(y_2-x_2)\bigr)\,
  F_{\rm out}(x_1,x_2)\,F_{\rm mid}(x_1,y_2)\,F_{\rm in}(y_1,y_2)
 \]
 has norm at most $r^2\norm{F_{\rm out}}_\infty\norm{F_{\rm mid}}_\infty
 \norm{F_{\rm in}}_\infty$, and so has the operator with $F_{\rm mid}(y_1,x_2)$ in
 place of $F_{\rm mid}(x_1,y_2)$.
\end{enumerate}
\end{lemma}

\begin{proof}
Let $C_1,\dots,C_{n_c}$ be the strongly connected components of $D_{\mathcal K}$,
ordered so that every arc between two components goes from an earlier to a later
one. Consider functions of mixed coordinates: $x_i$ for $i$ in the components
already treated and $y_i$ for the others, with uniform probability. Step $s$
replaces the coordinates $y_{C_s}$ by $x_{C_s}$:
\[
 g\longmapsto M_s\cdot B_sg,\qquad
 B_sg=\E_{y_{C_s}}\Bigl[\prod_{(i,l)\in\mathcal K\cap C_s^2}
 \chi\bigl(\varsigma_{il}(y_l-x_i)\bigr)\,g\Bigr],\qquad
 M_s=\prod_{(i,l)\in\mathcal K_s}\chi\bigl(\varsigma_{il}(y_l-x_i)\bigr),
\]
where $\mathcal K_s$ is the set of $(i,l)\in\mathcal K$ with $i\in C_s$ and $l$ in a
later component.
When step $s$ is performed the coordinates $x_{C_s}$ have just been created and
every $y_l$ with $l$ in a later component is still present, so $M_s$ is a function
of the current coordinates. Every pair of $\mathcal K$ is used exactly once: a
pair inside a component in its block $B_s$, and a pair between two components in
$M_s$ for the component $C_s$ of its first index, since its second index lies in a
later component. Hence $T_{\mathcal K}=M_{n_c}B_{n_c}\cdots M_1B_1$. Each $M_s$ is a
multiplication by a function of modulus at most one, and $B_s$ is the operator of
the block $C_s$ tensored with the identity in the other coordinates, which has the
same norm.

(a) If $D_{\mathcal K}$ is acyclic, every component is a single vertex $i$, and
$B_s$ is $h\mapsto\E_y\chi(\varsigma_{ii}(y-x))h(y)=\chi(\varsigma_{ii})Ch$ in the
coordinate $i$, with the Gauss operator $C$ of \cref{lem:spec:gauss}, of norm $r$.
So $\norm{T_{\mathcal K}}\le r^k$.

(b) For a two-vertex component $C_s=\{i,l\}$, strong connectivity gives both arcs,
so the kernel of $B_s$ is the complete rectangle
$\prod_{a,b\in\{i,l\}}\chi(\varsigma_{ab}(y_b-x_a))$. Let $Q_{\rm in}$ and
$Q_{\rm out}$ be the multiplications by $\one[y_i\ne y_l]$ and by
$\one[x_i\ne x_l]$. The steps before $s$ act on other coordinates or multiply
pointwise, so they commute with $Q_{\rm in}$, and $Q_{\rm in}P_{\rm in}=P_{\rm in}$.
The steps after $s$, and $M_s$, act on other coordinates or multiply pointwise, so
they commute with $Q_{\rm out}$, and $P_{\rm out}Q_{\rm out}=P_{\rm out}$. Hence
$Q_{\rm in}$ may be inserted immediately before $B_s$ and $Q_{\rm out}$ immediately
after it without changing $P_{\rm out}T_{\mathcal K}P_{\rm in}$, for every
two-vertex component. By \cref{cor:dl:normalized}(a) with $a=b=2$,
$\norm{Q_{\rm out}B_sQ_{\rm in}}\le2/p=2r^2$, and the singleton blocks have norm
$r$ as in (a). Multiplying, $\norm{P_{\rm out}T_{\mathcal K}P_{\rm in}}\le
r^{k-2\beta}(2r^2)^\beta=2^\beta r^k$.

(c) If $i_1\to i_2\to\dots\to i_\ell\to i_1$ is a directed cycle with distinct
vertices, the permutation with $\pi(i_s)=i_{s+1}$ (indices modulo $\ell$) and
$\pi(i)=i$ otherwise is a second perfect matching. Conversely, if $\pi\ne\mathrm{id}$
is a perfect matching in $\mathcal K$, every $i$ with $\pi(i)\ne i$ gives an arc
$i\to\pi(i)$, and following a cycle of $\pi$ gives a directed cycle.

(d) Let $M_F$ be the multiplication by $F$, and let $C_\varsigma\otimes I$ and
$I\otimes C_\varsigma$ be the Gauss operator
$(C_\varsigma h)(a)=\E_b\chi(\varsigma(b-a))h(b)$ of \cref{lem:spec:gauss} acting in the
first and in the second coordinate. Applied to $f(y_1,y_2)$, the operator multiplies
by $F_{\rm in}$, replaces $y_1$ by $x_1$ through the kernel
$\chi(\varsigma_{11}(y_1-x_1))$, multiplies by $F_{\rm mid}(x_1,y_2)$, replaces $y_2$
by $x_2$ through $\chi(\varsigma_{22}(y_2-x_2))$, and multiplies by $F_{\rm out}$: it
is $M_{F_{\rm out}}(I\otimes C_{\varsigma_{22}})M_{F_{\rm mid}}(C_{\varsigma_{11}}\otimes
I)M_{F_{\rm in}}$. With $F_{\rm mid}(y_1,x_2)$ the two transforms are applied in the
other order. Each Gauss operator has norm $r$, and each multiplication has norm at
most the supremum of its function.
\end{proof}

\begin{remark}[The witness kernels that occur]\label{rem:extraction-kernels}
For a witness of size two the rectangle is a matching, a three-edge path (one
connector present) or a complete rectangle (both present). The first two are
acyclic, so their kernels have norm at most $r^2=1/p$ for independent uniform
labels, with no injectivity condition (the \emph{incomplete-rectangle bound}). The
complete rectangle has norm at most $2r^2$ on distinct pairs; numerical evidence
for the sharpness of $2$ is discussed in \cref{rem:dl:K-sharp}. A witness of size
three or four found along an augmenting path is used only if its digraph has no
 strongly connected component of three or more vertices in the selected
 partition. A complete three-vertex component may instead be priced at
 $6r^3$ by \cref{thm:cubic-resultant}, with internal distinctness on
 both sides. This extension does not apply to an incomplete three-vertex
 component. The original criterion holds, for instance,
when the connectors between a new witness pair and the others all point one way:
the new vertex then lies on no directed cycle, and the kernel factors through one
more Gauss operator. Part~(d) keeps inside the norm a bounded multiplier that depends
on one endpoint of each of the two pairs of a matching; a third character
$\chi(\varsigma_{12}(y_2-x_1))$ or $\chi(\varsigma_{21}(y_1-x_2))$ has modulus at most
one and may be absorbed into $F_{\rm mid}$, so the three-edge paths are covered too.
In the applications the other coordinates of a larger witness are fixed first and the
multipliers may depend on them; the score multipliers of the positive roots are of
this form (\cref{sec:positive-roots}). No Schur-multiplier bound for a signed kernel
is used.
\end{remark}

\subsection{A bound for each class $\mathcal C(P,A)$}\label{subsec:extraction:extraction}

A class is bounded in four steps. (1)~Its character sum is written
as the witness kernel times the nonnegative factors $1+\chi_e$ of
the free pairs. (2)~One conditions on the labels of a vertex
cover. (3)~The witness labels then have a bounded density, and
their kernel is bounded by \cref{lem:matching:kernel}.
(4)~The remaining labels are integrated one vertex at a time,
each contributing its vertex factor $F_v$ of
\cref{def:extraction:records}, with an
additional second-order term when a label coincides with a fixed
label of the other row. The definitions below record exactly the
choices (witness, cover, orders) and the factor contributed by
each non-witness vertex.

\begin{definition}[Records of a class (extraction data)]\label{def:extraction:records}
Let $\mathcal C=\mathcal C(P,A)$ be a class with free pairs $F$, a witness of
size $k$ with sets $V\subseteq X$, $W\subseteq Z$, and a cover $(C_x,C_z)$ with
integrated sets $U_x$, $U_z$ (\cref{def:extraction:witness}). Fix orders of the
four sets $U_x$, $U_z$, $C_x$, $C_z$. Witness, cover and orders form a
\emph{record} of $\mathcal C$ (extraction data, not a record in the sense of
\cref{subsec:intro:glossary}). For a \emph{non-witness vertex}
$v\in(X\setminus V)\cup(Z\setminus W)$ define $t_v$, the set $H_v$ of its \emph{free
partners}, $h_v=|H_v|$ and $s_v=t_v+h_v$:
\begin{itemize}
\item $u\in U_x$: $t_u$ is the number of $G^+$-neighbours of $u$ in $V\cup C_x$
 or earlier in $U_x$, and $H_u=\{b:(u,b)\in F\}$;
\item $u\in U_z$: $t_u$ is the number of $G^+$-neighbours of $u$ in $W\cup C_z$
 or earlier in $U_z$, and $H_u=\{a:(a,u)\in F\}$;
\item $b\in C_z$: $t_b$ is the number of $G^+$-neighbours of $b$ earlier in
 $C_z$, and $H_b=\{a:(a,b)\in F\}$ (all free pairs at $b$, also those from
 $V$);
\item $a\in C_x$: $t_a$ is the number of $G^+$-neighbours of $a$ earlier in
 $C_x$, and $H_a=\{b:(a,b)\in F\}$.
\end{itemize}
Viewed in $U^+$, $H_v\subseteq N(v)$, by the symmetry of $N$; and by \textup{(VC$_k$)},
$H_u\subseteq C_z$ for $u\in U_x$ and $H_u\subseteq C_x$ for $u\in U_z$. Thus
$t_v\le d^0_v$, $h_v\le r_v$ and $s_v\le d_v\le\Delta$. Put
$t_V=e(V)+e(V,C_x)$ and $t_W=e(W)+e(W,C_z)$, let $f_V$ and $f_W$ be the numbers
of free pairs with $x$-endpoint in $V$ and with $x'$-endpoint in $W$, and let
$\mathrm{ex}(\mathcal C)=t_V+t_W+f_V+f_W$ be the \emph{exponent} of the record.
With the \emph{merge counts} (for the merged roots of
\cref{lem:extraction-merged-count} below, not the merges of
\cref{sec:completion-repair})
\[
 \mathrm{mg}(j,i)=\#\{i'\in\operatorname{Adj}(i):j\in N(i')\},\qquad
 c^U(j,i)=\#\{u\in\operatorname{Adj}(j):u\in N(i)\}\qquad(i,j\in U^+),
\]
the \emph{vertex factors} of the record are, with Iverson brackets,
\begin{align*}
 F_u&=\frac{1+e_{s_u}(p,n)+r^22^{s_u}\sum_{b\in C_z}2^{c^U(u,b)-[b\in H_u]}}
 {1-e_{t_u}(p,n)}&&(u\in U_x),\\
 F_u&=\frac{1+e_{s_u}(p,n)+r^22^{s_u}\sum_{a\in C_x}2^{c^U(u,a)-[a\in H_u]}}
 {1-e_{t_u}(p,n)}&&(u\in U_z),\\
 F_b&=\frac{1+e_{s_b}(p,2n)+r^22^{s_b}\sum_{j\in U^+}
 2^{\mathrm{mg}(j,b)+c^U(j,b)-[j\in H_b]}}{1-e_{t_b}(p,n)}&&(b\in C_x\cup C_z).
\end{align*}
The brackets record that the free factor of a pair whose two labels coincide is
$1+\chi(0)=1$. There are exactly $2(n-k)$ non-witness vertices, hence $2(n-k)$
vertex factors. A class may carry several records; each choice is made once,
independently of $p$.
\end{definition}

\begin{lemma}[Counts with merged roots]\label{lem:extraction-merged-count}
Let $\beta_1,\dots,\beta_t\in\Fp$ be distinct, let $\gamma_1,\dots,\gamma_h\in\Fp$
be distinct, let $s_1,\dots,s_t,s'_1,\dots,s'_h\in\{\pm1\}$, and let
$\mathcal X\subseteq\Fp$ contain all $\beta_i$ and $\gamma_l$. Let
$\mu_{\rm m}=\#\{(i,l):\beta_i=\gamma_l\}$, the number of \emph{merged} roots. Then
for every $n'\ge0$
\[
 \sum_{\substack{y\in\Fp\setminus\mathcal X\\ \chi(y-\beta_i)=s_i\ \forall i}}
 \ \prod_{l=1}^h\bigl(1+s'_l\chi(y-\gamma_l)\bigr)
 \ \le\ p\,2^{-t}2^{\mu_{\rm m}}\bigl(1+e_{t+h-\mu_{\rm m}}(p,n')\bigr)
 \ \le\ p\,2^{-t}2^{\mu_{\rm m}}\bigl(1+e_{t+h}(p,n')\bigr).
\]
\end{lemma}

\begin{proof}
For $y\notin\mathcal X$ every factor $1+s'_l\chi(y-\gamma_l)$ is $0$ or $2$, so the
sum is $2^h$ times the number of $y\notin\mathcal X$ satisfying the $t+h$
conditions $\chi(y-\beta_i)=s_i$ and $\chi(y-\gamma_l)=s'_l$. Since the $\beta_i$
are distinct and the $\gamma_l$ are distinct, the merged roots pair one $\beta_i$
with one $\gamma_l$. Two conditions at the same root are identical or
incompatible; in the second case the sum is $0$. Otherwise the conditions are
$t+h-\mu_{\rm m}$ conditions with distinct roots, all in $\mathcal X$, and by
\cref{lem:extension}\ref{it:ext:upper} and \ref{it:et:compare} their solution set
has at most $p2^{-(t+h-\mu_{\rm m})}(1+e_{t+h-\mu_{\rm m}}(p,n'))$ elements.
Multiply by $2^h$; the second inequality holds because $e_t(p,n')$ is
nondecreasing in $t$ \ref{it:et:poly}.
\end{proof}

\begin{lemma}[Cover rows]\label{lem:extraction-rows}
Assume \textup{(S1)} at $p$. Let $C\subseteq U^+$ be a set of $x'$-vertices with an
order, let $x\in\mathcal I^+$ and $H\subseteq\{(a,b)\in E_\times:b\in C\}$. For
$b\in C$ let $t_b$ be the number of $G^+$-neighbours of $b$ earlier in $C$,
$H_b=\{a:(a,b)\in H\}$, $h_b=|H_b|$, $s_b=t_b+h_b$, and let $F_b(H)$ be the
expression of $F_b$ in \cref{def:extraction:records} with these counts and these
sets $H_b$. Let $\mu_C$ be the law of
$x'_C$ for $x'\sim\mu$, and
$\kappa_C(x,c')=\sum_{b\in C}\sum_{i\in U^+}c^U(i,b)\one[c'_b=x_i]$. Then
\[
 \E_{c'\sim\mu_C}\Bigl[\prod_{(a,b)\in H}\bigl(1+\chi(\varsigma_{ab}(c'_b-x_a))\bigr)
 \,2^{\kappa_C(x,c')}\Bigr]\ \le\ R_{U^+\setminus C}\prod_{b\in C}F_b(H).
\]
\end{lemma}

\begin{proof}
Call the integrand $\Pi(c')\ge0$.

\emph{Step 1 (reference law).} Let $\lambda_C=\lambda_p(G^+[C])$. By
\cref{lem:graphs:vertex}\ref{it:vertex:subset} the density of $\mu_C$ with respect
to $\lambda_C$ takes values in an interval $[c_0,R_{U^+\setminus C}c_0]$; since it has
mean one, $c_0\le1$ and the density is at most $R_{U^+\setminus C}$. So it suffices
to show $\E_{\lambda_C}\Pi\le\prod_{b\in C}F_b(H)$.

\emph{Step 2 (lower count).} By \cref{lem:graphs:sequential}\ref{it:seq:count}
with the error pair of the setting, no anchors and the given order,
\[
 |\operatorname{Inj}_p(G^+[C])|\ge\prod_{b\in C}p2^{-t_b}\bigl(1-e_{t_b}(p,n)\bigr).
\]
Write $\Pi(c')=\prod_b\omega_b(c'_b)$ with
\[
 \omega_b(y)=\prod_{a\in H_b}\bigl(1+\chi(\varsigma_{ab}(y-x_a))\bigr)\cdot
 2^{\sum_ic^U(i,b)\one[y=x_i]}.
\]

\emph{Step 3 (merge weights).} Call $b$ \emph{coincident} if $c'_b=x_{j(b)}$ for
some (unique) $j(b)\in U^+$. A \emph{merge at $b$} is a pair $(b',a)$ with $b'$
earlier than $b$, $b'\sim b$, $a\in H_b$ and $c'_{b'}=x_a$; let
$\mu_b=\mu_b(c'_{<b})$ be their number. If $(b',a)$ is a merge at $b$, then
$b'$ is coincident with $j(b')=a$, and $b\in\operatorname{Adj}(b')$ with
$j(b')\in N(b)$. So a fixed $b'$ causes at most $\mathrm{mg}(j(b'),b')$
merges, and $\sum_b\mu_b\le\sum_{b'\ \text{coincident}}\mathrm{mg}(j(b'),b')$. With $\hat\omega_b(y)=\omega_b(y)2^{\mathrm{mg}(j,b)}$ if
$y=x_j$ for some $j\in U^+$, and $\hat\omega_b(y)=\omega_b(y)$ otherwise,
\[
 \Pi(c')\le\prod_b\hat\omega_b(c'_b)\,2^{-\mu_b(c'_{<b})}.
\]

\emph{Step 4 (one vertex).} Fix $b$ and the earlier labels $c'_{<b}$, and let
$\mathrm{Adm}_b$ be the set of $y$ distinct from the earlier labels that satisfy the
$t_b$ constraints of $G^+$ with the earlier neighbours of $b$. We claim
\begin{equation}\label{eq:extraction:onevertex}
 \sum_{y\in\mathrm{Adm}_b}\hat\omega_b(y)\,2^{-\mu_b}\le p2^{-t_b}
 \Bigl(1+e_{s_b}(p,2n)+r^22^{s_b}\sum_{j\in U^+}
 2^{\mathrm{mg}(j,b)+c^U(j,b)-[j\in H_b]}\Bigr).
\end{equation}
For $y\notin x(U^+)$, $\hat\omega_b(y)=\prod_{a\in H_b}(1+\chi(\varsigma_{ab}(y-x_a)))$.
The internal roots (the labels of the earlier neighbours) are distinct, the cross
roots $x_a$ ($a\in H_b$) are distinct, and they coincide exactly at the $\mu_b$
merges. With $\mathcal X=x(U^+)\cup\{c'_{b'}:b'\text{ earlier}\}$,
\cref{lem:extraction-merged-count} bounds the sum of $\hat\omega_b(y)$ over these
values by $p2^{-t_b}2^{\mu_b}(1+e_{s_b}(p,2n))$, so their contribution to the left
side of \eqref{eq:extraction:onevertex} is at most $p2^{-t_b}(1+e_{s_b}(p,2n))$. For $y=x_j$ the factor of $a=j$, if
$j\in H_b$, is $1+\chi(0)=1$, the other factors are at most $2$, and $x$ is
injective; so
$\hat\omega_b(x_j)2^{-\mu_b}\le2^{h_b-[j\in H_b]+c^U(j,b)+\mathrm{mg}(j,b)}
=p2^{-t_b}r^22^{s_b+\mathrm{mg}(j,b)+c^U(j,b)-[j\in H_b]}$, because $pr^2=1$.
Summing over $j\in U^+$ proves \eqref{eq:extraction:onevertex}.

\emph{Step 5 (product).} The right side of \eqref{eq:extraction:onevertex} does not
depend on $c'_{<b}$, and $\mu_b$ depends only on $c'_{<b}$. Summing the bound of
Step 3 over $c'\in\operatorname{Inj}_p(G^+[C])$, successively over the labels in reverse order and using the uniform
bound \eqref{eq:extraction:onevertex} for each label given the earlier ones,
gives $\sum_{c'}\Pi(c')\le\prod_bp2^{-t_b}(1+e_{s_b}(p,2n)+r^22^{s_b}\sum_j
2^{\mathrm{mg}(j,b)+c^U(j,b)-[j\in H_b]})$, and dividing by the count of Step 2
gives $\E_{\lambda_C}\Pi\le\prod_bF_b(H)$.
\end{proof}

\begin{theorem}[Grouped extraction with a witness of size $k$]\label{lem:matching:extractionk}
Assume \textup{(S1)} at $p$. Let $\mathcal C=\mathcal C(P,A)$ be a class with a
record (\cref{def:extraction:records}) whose witness has size $k$, and suppose that
the kernel
\[
 \mathsf W(v,w)=\prod_{(a,b)\in P\cap(V\times W)}\chi\bigl(\varsigma_{ab}(w_b-v_a)\bigr)
 \qquad(v\in\Fp^V,\ w\in\Fp^W)
\]
of its present rectangle pairs, compressed to injective tuples, has norm at most
$c_{\rm ker}r^k$ for uniform probability; by \cref{lem:matching:kernel},
$c_{\rm ker}=2^{\beta(\mathcal C)}$ when every strongly connected component of
$D_{\mathcal C}$ has at most two vertices. For either integrated set put
$\widetilde R_U=\prod_{u\in U}(1+e_{t_u}(p,n))/(1-e_{t_u}(p,n))$,
where $t_u$ is its actual internal insertion degree. Then for all $f,h\in L^2(\mu)$
\[
 \Bigl|\sum_{S\in\mathcal C}\ip f{T_Sh}_\mu\Bigr|
 \le c_{\rm ker}\,r^k\,2^{\mathrm{ex}(\mathcal C)/2}\,
 \frac{\sqrt{\widetilde R_{U_x}\widetilde R_{U_z}}}{(1-\hat e)^k}
 \prod_{v\in(X\setminus V)\cup(Z\setminus W)}F_v^{1/2}\;\norm f_\mu\norm h_\mu .
\]
\end{theorem}

\begin{proof}
Let $E_{VW}$ be the rectangle and $F$ the free pairs.

\emph{Step 1 (grouping).} The sets of $\mathcal C$ are the sets $P\cup S'$ with
$S'\subseteq F$, so
$\sum_{S\in\mathcal C}\chi_S=\prod_{e\in P}\chi_e\prod_{e\in F}(1+\chi_e)
=\prod_{e\in P\cup F}\psi_e$, with $\psi_e=\chi_e$, $|\psi_e|\le1$, for $e\in P$ and
$\psi_e=1+\chi_e\in[0,2]$ for $e\in F$. By (W$_k$) no pair of $E_{VW}$ is free. Let
$(a,b)$ be a non-absent pair outside $E_{VW}$. If $a\in V\cup U_x$, then
$a\notin C_x$ and $(a,b)\notin V\times W$, so $b\in C_z$ by (VC$_k$). Hence every
such pair lies in exactly one of the groups
\[
 \mathcal G_1:\ a\in V\cup U_x,\ b\in C_z;\qquad
 \mathcal G_2:\ a\in C_x,\ b\in W\cup U_z;\qquad
 \mathcal G_\Gamma:\ a\in C_x,\ b\in C_z,
\]
and $\sum_{S\in\mathcal C}\chi_S=\mathsf W(v,w)\,\Theta_1\Theta_2\Gamma$ with
$v=x_V$, $w=x'_W$ and $\Theta_1,\Theta_2,\Gamma$ the products of $\psi_e$ over
$\mathcal G_1,\mathcal G_2,\mathcal G_\Gamma$. Here $\Theta_1$ depends on
$x_{V\cup U_x}$ and $x'_{C_z}$, $\Theta_2$ on $x_{C_x}$ and $x'_{W\cup U_z}$, and
$\Gamma$ on $x_{C_x}$ and $x'_{C_z}$. Let $\Theta^+_1$ and $\Gamma^+$ be the
products of $1+\chi_e$ over the free pairs of $\mathcal G_1$ and of
$\mathcal G_\Gamma$; then $|\Theta_1|\le\Theta^+_1$ and $|\Gamma|\le\Gamma^+$. A free
pair with $x$-endpoint in $V$ lies in $\mathcal G_1$, so
$\Theta^+_1=\Theta^+_V\Theta^+_U$ with
$\Theta^+_V=\prod_{(a,b)\in F,\,a\in V}(1+\chi_{ab})\le2^{f_V}$ and
$\Theta^+_U=\prod_{(u,b)\in F,\,u\in U_x}(1+\chi_{ub})$.

\emph{Step 2 (conditioning).} Put $c=x_{C_x}$ and $c'=x'_{C_z}$. Under
$\mu\otimes\mu$ the labellings $x$ and $x'$ are independent, so given $(c,c')$ the
pairs $(x_V,x_{U_x})$ and $(x'_W,x'_{U_z})$ are independent. Let $\nu_V$ and
$\nu_W$ be the conditional laws of $x_V$ given $c$ and of $x'_W$ given $c'$ under
$\mu$, and put
$\mathsf a(v)=\E[f(x)\Theta_1\mid x_V=v,c,c']$ and
$\mathsf b(w)=\E[\Theta_2h(x')\mid x'_W=w,c,c']$. The characters being real,
\[
 \sum_{S\in\mathcal C}\ip f{T_Sh}_\mu
 =\E_{\mu\otimes\mu}\Bigl[\overline{f(x)}\sum_{S\in\mathcal C}\chi_S\,h(x')\Bigr]
 =\E_{c,c'}\Bigl[\Gamma\,\E_{v\sim\nu_V,\,w\sim\nu_W}
 \bigl[\overline{\mathsf a(v)}\,\mathsf W(v,w)\,\mathsf b(w)\bigr]\Bigr].
\]

\emph{Step 3 (density together with its cover marginal).}
Keep the actual conditional density supremum
$D_V(c)=\sup_v\tilde\nu_V(v\mid c)$ as a function of the cover labels.
Write $N_C(c)$ for the number of completions on $V\sqcup U_x$,
$N_U(v,c)$ for the number on $U_x$, and $\lambda_C$ for the uniform
strict labelling law on $C_x$. Exact counting gives
\[
 \frac{D_V(c)\mu_{C_x}(c)}{\lambda_C(c)}
 =\frac{p^k|\mathcal I_{C_x}|\max_v N_U(v,c)}{|\mathcal I^+|}.
\]
All cover labellings extend by (S1). Inserting $V$ after $C_x$ gives
at least $p^k2^{-t_V}(1-\hat e)^k$ choices. The subsequent insertion
of $U_x$ gives lower and upper products
$\prod_up2^{-t_u}(1-e_{t_u})$ and
$\prod_up2^{-t_u}(1+e_{t_u})$. Thus, as measures,
\begin{equation}\label{eq:extraction:cross-density}
 D_V\,d\mu_{C_x}\le
 \frac{2^{t_V}\widetilde R_{U_x}}{(1-\hat e)^k}\,d\lambda_{C_x}.
\end{equation}
The same assertion holds for $D_W(c')$ and $C_z$.

\emph{Step 4 (the witness kernel).} The functions $\mathsf a\tilde\nu_V$ and
$\mathsf b\tilde\nu_W$ vanish off injective tuples, so the kernel hypothesis and
$\E|\mathsf a|^2\tilde\nu_V^2\le D_V(c)\E|\mathsf a|^2\tilde\nu_V$ (expectations for
uniform probability) give
$|\E_{\nu_V\otimes\nu_W}[\overline{\mathsf a}\mathsf W\mathsf b]|
=|\ip{\mathsf a\tilde\nu_V}{T_{\mathsf W}(\mathsf b\tilde\nu_W)}|
\le c_{\rm ker}r^k\sqrt{D_V(c)D_W(c')}\norm{\mathsf a}_{\nu_V}\norm{\mathsf b}_{\nu_W}$.
Hence
\begin{equation}\label{eq:extraction:step4}
 \Bigl|\sum_{S\in\mathcal C}\ip f{T_Sh}_\mu\Bigr|\le c_{\rm ker}r^k
 \E_{c,c'}\bigl[|\Gamma|\sqrt{D_V(c)D_W(c')}\,
       \norm{\mathsf a}_{\nu_V}\norm{\mathsf b}_{\nu_W}\bigr].
\end{equation}

\emph{Step 5 (Cauchy--Schwarz in the integrated variables).} With the weight
$|\Theta_1|\le\Theta^+_1$,
$|\mathsf a(v)|^2\le\E[|f|^2\Theta^+_1\mid v,c,c']\,\E[\Theta^+_1\mid v,c,c']$. Given
$(v,c)$ the labels $x_{U_x}$ are uniform on the set $\mathcal Y(v,c)$ of completions
of $(v,c)$ to an element of $\mathcal I^+$, so
\[
 \E[\Theta^+_1\mid v,c,c']=\Theta^+_V\,\mathsf m(v,c,c')\le2^{f_V}\mathsf m(v,c,c'),
 \qquad\mathsf m(v,c,c')=\frac1{|\mathcal Y|}\sum_{y\in\mathcal Y}\Theta^+_U(y,c').
\]

\emph{Step 6 (the integrated vertices).} Label $U_x$ in its order, writing $y_u$ for
the label of $u$. For $u\in U_x$ the free partners are $H_u=\{b:(u,b)\in F\}$; a free
pair $(u,b)$ is not in $E_{VW}$ and $u\notin C_x$, so $H_u\subseteq C_z$ by (VC$_k$),
and
$\Theta^+_U=\prod_uw_u(y_u)$ with
$w_u(y)=\prod_{b\in H_u}(1+\chi(\varsigma_{ub}(c'_b-y)))$. A \emph{merge at $u$} is a
pair $(i,b)$ with $i\in\operatorname{Adj}(u)$ labelled before $u$ (so $i\in V\cup
C_x$, or $i\in U_x$ earlier than $u$), $b\in H_u$, and label of $i$ equal to
$c'_b$. Let $\kappa^f_u$ count the merges at $u$ with $i\in V\cup C_x$, which are
determined by $(v,c,c')$, let $\kappa^r_u(y_{<u})$ count those with $i\in U_x$, and
$\kappa^f=\sum_u\kappa^f_u$. Put $\hat w_u(y)=w_u(y)2^{c^U(u,b)}$ if $y=c'_b$ for
some $b\in C_z$ (unique, as $c'$ is injective) and $\hat w_u(y)=w_u(y)$ otherwise.
If $(i,b)$ is a merge at $u$ with $i\in U_x$, then $y_i=c'_b$,
$u\in\operatorname{Adj}(i)$ and $u\in N(b)$; so an $i\in U_x$ with $y_i=c'_b$
causes at most $c^U(i,b)$ merges, and
$\prod_uw_u(y_u)\le\prod_u\hat w_u(y_u)2^{-\kappa^r_u(y_{<u})}$.

Fix $u$ and the earlier labels, and sum over the admissible values $y$ of $y_u$
(distinct from the labels already used, satisfying the $t_u$ constraints with the
labelled neighbours of $u$). For the \emph{generic values}
$y\notin\{c'_b:b\in C_z\}$, $\hat w_u(y)=w_u(y)$; the internal roots (labels of
the labelled neighbours) and the cross roots $c'_b$ ($b\in H_u$) merge exactly at
the $\kappa^f_u+\kappa^r_u$ merges at $u$. With $\mathcal X$ the set of labels
already used together with $\{c'_b:b\in C_z\}$, \cref{lem:extraction-merged-count}
bounds the sum of $\hat w_u(y)$ over these values by
$p2^{-t_u}2^{\kappa^f_u+\kappa^r_u}(1+e_{s_u}(p,n))$.
For the \emph{special values} $y=c'_b$, $b\in C_z$, every factor of $w_u(c'_b)$ is at
most $2$, and that of $b$ itself, if $b\in H_u$, is $1+\chi(0)=1$; so
$\hat w_u(c'_b)\le2^{h_u-[b\in H_u]+c^U(u,b)}=p2^{-t_u}r^22^{s_u+c^U(u,b)-[b\in H_u]}$;
distinct $b$ give distinct values. Hence, uniformly in the earlier labels,
\[
 \sum_y\hat w_u(y)\,2^{-\kappa^r_u}\le p2^{-t_u}2^{\kappa^f_u}
 \Bigl(1+e_{s_u}(p,n)+r^22^{s_u}\sum_{b\in C_z}2^{c^U(u,b)-[b\in H_u]}\Bigr).
\]
Summing over $y\in\mathcal Y$ with the last vertex first bounds
$\sum_{y\in\mathcal Y}\Theta^+_U(y,c')$ by the product of these bounds. Since $(v,c)$
is feasible, \cref{lem:graphs:sequential}\ref{it:seq:count} with the error pair of the
setting, the anchors $V\cup C_x$ and the order of $U_x$ gives
$|\mathcal Y|\ge\prod_{u\in U_x}p2^{-t_u}(1-e_{t_u}(p,n))$. Dividing,
\begin{equation}\label{eq:extraction:mbound}
 \mathsf m(v,c,c')\le2^{\kappa^f}\prod_{u\in U_x}F_u .
\end{equation}
Finally $\kappa^f\le\kappa_{C_z}(x,c')$ for every completion $x$ of $(v,c)$, with
$\kappa_{C_z}$ of \cref{lem:extraction-rows}: a pair $(i,b)$ with $x_i=c'_b$ is a
merge at most at the $c^U(i,b)$ vertices of $\operatorname{Adj}(i)\cap N(b)$.

\emph{Step 7 (the cover rows).} Apply Cauchy--Schwarz to
\eqref{eq:extraction:step4}, placing $D_W(c')$ with
$\norm{\mathsf a}_{\nu_V}^2$ and $D_V(c)$ with
$\norm{\mathsf b}_{\nu_W}^2$. By Steps 5 and 6, since $\Gamma^+$ depends only on
$(c,c')$,
$|\Gamma|\,|\mathsf a(v)|^2\le2^{f_V}\prod_{u\in U_x}F_u\,
\E[|f(x)|^2\Theta^+_1\Gamma^+2^{\kappa_{C_z}(x,c')}\mid v,c,c']$. Integrating over
$v\sim\nu_V$ and over $(c,c')\sim\mu_{C_x}\otimes\mu_{C_z}$, under which $x\sim\mu$
and $c'\sim\mu_{C_z}$ are independent,
\[
 \E_{c,c'}\bigl[|\Gamma|D_W(c')\norm{\mathsf a}^2_{\nu_V}\bigr]\le2^{f_V}\prod_{u\in U_x}F_u\;
 \E_{x\sim\mu}\Bigl[|f(x)|^2\,\E_{c'\sim\mu_{C_z}}
 \bigl[D_W(c')\Theta^+_1\Gamma^+2^{\kappa_{C_z}(x,c')}\bigr]\Bigr].
\]
The product $\Theta^+_1\Gamma^+=\Theta^+_V\Theta^+_U\Gamma^+$ runs over exactly the
free pairs with $x'$-endpoint in $C_z$, their $x$-endpoints lying in $V$, $U_x$ or
$C_x$. \Cref{lem:extraction-rows} with $C=C_z$ in its order and
$H=\{(a,b)\in F:b\in C_z\}$, whose counts and sets $H_b$ are those of
\cref{def:extraction:records}, is applied after
\eqref{eq:extraction:cross-density} has replaced
$D_W\,d\mu_{C_z}$ by its multiple of $d\lambda_{C_z}$.
Only Steps 2--5 of that row lemma are then needed; its separate
marginal comparison is omitted. The inner expectation is at most
$2^{t_W}\widetilde R_{U_z}(1-\hat e)^{-k}\prod_{b\in C_z}F_b$.
Consequently
\[
 \E_{c,c'}[|\Gamma|D_W\norm{\mathsf a}^2_{\nu_V}]
 \le\frac{2^{f_V+t_W}\widetilde R_{U_z}}{(1-\hat e)^k}
 \prod_{u\in U_x}F_u\prod_{b\in C_z}F_b\,\norm f_\mu^2.
\]

\emph{Step 8 (the other side, and conclusion).} The map $(x,x')\mapsto(x',x)$
preserves $\mu\otimes\mu$, and
$\chi_{(a,b)}(x,x')=\chi(\varsigma_{ab}(x'_b-x_a))=\chi(\varsigma_{ba}(x_a-x'_b))
=\chi_{(b,a)}(x',x)$ by \cref{lem:spec:elementary}. With the two labellings
exchanged, $\mathsf b$ is therefore the function $\mathsf a$ of the transposed class
$\mathcal C(P^{\mathsf T},A^{\mathsf T})$, $P^{\mathsf T}=\{(b,a):(a,b)\in P\}$, of the
same spectral data, with $h$ in place of $f$; its record is the transposed
witness with the same cover and orders, $(W,C_z,U_z)$ playing the role of
$(V,C_x,U_x)$. Since both labellings carry the half graph $G^+$ and $N$ is
symmetric, the counts and factors of the transposed record are those defined in
\cref{def:extraction:records} for $U_z$ and $C_x$. Steps 5--7 for the transposed
class give
$\E_{c,c'}[|\Gamma|D_V\norm{\mathsf b}^2_{\nu_W}]$ bounded by
$2^{f_W+t_V}\widetilde R_{U_x}(1-\hat e)^{-k}
\prod_{u\in U_z}F_u\prod_{a\in C_x}F_a\,\norm h_\mu^2$.
The cross-weighted Cauchy--Schwarz inequality of Step 7 gives
\[
 \Bigl|\sum_{S\in\mathcal C}\ip f{T_Sh}_\mu\Bigr|
 \le c_{\rm ker}r^k\,2^{\mathrm{ex}(\mathcal C)/2}
 \frac{\sqrt{\widetilde R_{U_x}\widetilde R_{U_z}}}{(1-\hat e)^k}
 \prod_{v\in(X\setminus V)\sqcup(Z\setminus W)}F_v^{1/2}
 \norm f_\mu\norm h_\mu .
\]
Here $\mathrm{ex}=t_V+t_W+f_V+f_W$. The vertex factors left are those of
$U_x\sqcup C_z\sqcup U_z\sqcup C_x=(X\setminus V)\sqcup(Z\setminus W)$. This is the
claim. Since each insertion degree is at most its full internal degree,
$\sqrt{\widetilde R_{U_x}\widetilde R_{U_z}}\le R_{\rm full}$.
In particular the earlier, coarser factor $R_{\rm full}^2$ remains valid
for the retained leaf and histogram certificates below.
\end{proof}

\begin{remark}[What the proof uses]\label{rem:extraction-proof}
(1) Only Steps 3 and 4 see the size $k$ of the witness, through the density of the
witness labels and the norm of the witness kernel; the other steps use only the
cover condition and the counts $s_v=t_v+h_v$. So witnesses of size three and four
need no hypothesis beyond \cref{lem:matching:kernel}.
(2) The cover may be chosen for each class. With $C_x=X\setminus V$ and
$C_z=Z\setminus W$ no vertex is integrated. Integrated vertices may carry free
pairs; only their coincidences with fixed labels of the other row produce the weight
$2^{\kappa^f}$, which the cover rows absorb at the second-order cost
$r^22^{s_b}\sum_j2^{\mathrm{mg}(j,b)+c^U(j,b)-[j\in H_b]}$. Each merge costs one factor $2$,
moved to the coincident vertex that caused it, so the bound factorizes over the
coincidence patterns.
(3) The count $h_b$ of a cover vertex must include its free pairs to the witness set
$V$: their factors occur in Step 7 a second time, inside the product integrated by
\cref{lem:extraction-rows}, after the first copy was bounded by $2^{f_V}$ in Step 5.
A vertex $u\in U_x$ has no free pair to $W$, by (VC$_k$), so $h_u$ is its full free
degree.
(4) Step 8 is the only use of the symmetry of $N$. Without it the transposed class
has its own counts, taken in the column neighbourhoods
(\cref{subsec:push:pinnedremainder}).
\end{remark}

The certificates evaluate two consequences of
\cref{lem:matching:extractionk}. The first keeps the leaf-specific counts $s_v$,
bounds each special sum by its largest value over the vertices and the free
partners that the count $s_v$ allows, and bounds the denominators uniformly; the
second averages the labelling orders and is linear in the logarithm. For
$0\le s\le\Delta$ and $i\in U^+$ put $h^0_i(s)=\max(0,s-d^0_i)$,
$w^{\rm m}_{ij}=2^{c^U(i,j)}$ and $w^\varrho_{ij}=2^{\mathrm{mg}(j,i)+c^U(j,i)}$
($j\in U^+$), and let $\sigma^{\rm m}_i(s)$ and $\sigma^\varrho_i(s)$ be the sums of
the $h^0_i(s)$ smallest of the numbers $w^{\rm m}_{ij}$, respectively
$w^\varrho_{ij}$, over $j\in N(i)$; if $d_i\ge s$, then $h^0_i(s)\le r_i$. Put
\begin{equation}\label{eq:leaf-factors}
\begin{gathered}
 \mathsf S_{\rm m}(s)=\max_{i\in U^+,\ d_i\ge s}\Bigl(\sum_{j\in U^+}w^{\rm m}_{ij}
 -\tfrac12\sigma^{\rm m}_i(s)\Bigr),\\
 \mathsf S_\varrho(s)=\max_{i\in U^+,\ d_i\ge s}\Bigl(\sum_{j\in U^+}w^\varrho_{ij}
 -\tfrac12\sigma^\varrho_i(s)\Bigr),\\
 \mathsf f_{\rm m}(s)=\frac{1+e_s(p,n)+r^22^s\mathsf S_{\rm m}(s)}{1-\hat e},\qquad
 \mathsf f_\varrho(s)=\frac{1+e_s(p,2n)+r^22^s\mathsf S_\varrho(s)}{1-\hat e}.
\end{gathered}
\end{equation}
The maxima are over nonempty sets, since $\Delta=\max_id_i$. The constants
$\mathsf S_{\rm m}(s)$ and $\mathsf S_\varrho(s)$ are at most
$\max_i\sum_jw^{\rm m}_{ij}$ and $\max_i\sum_jw^\varrho_{ij}$, the values that bound
the special sums when every factor at a special value is bounded by $2$. The
subtracted terms are exact maxima: for a set $H\subseteq N(i)$ of at least
$h^0_i(s)$ free partners,
$\sum_jw_{ij}2^{-[j\in H]}=\sum_jw_{ij}-\frac12\sum_{j\in H}w_{ij}$, which is
largest when $H$ consists of $h^0_i(s)$ elements of smallest weight.

\begin{corollary}[Leaf bound and nonstar coefficients]\label{cor:extraction-leaf}
\leavevmode
\begin{enumerate}[label=\textup{(\alph*)}]
\item In \cref{lem:matching:extractionk} with $c_{\rm ker}=2^{\beta(\mathcal C)}$ put
 $\mathrm{ex}^\star(\mathcal C)=\mathrm{ex}(\mathcal C)+2\beta(\mathcal C)$. Then
 \[
  \norm{T_{\mathcal C}}\le r^k\,2^{\mathrm{ex}^\star(\mathcal C)/2}\,
  \frac{R_{\rm full}^2}{(1-\hat e)^k}\prod_{u\in U_x\cup U_z}\mathsf f_{\rm m}(s_u)^{1/2}
  \prod_{b\in C_x\cup C_z}\mathsf f_\varrho(s_b)^{1/2}.
 \]
 If moreover the \emph{monotone-factor condition} (on the
 factor functions $\mathsf S_{\rm m},\mathsf S_\varrho$ of
 \eqref{eq:leaf-factors})
 \begin{equation}\label{eq:monotone-signature}
  2^{s+1}\mathsf S_{\rm m}(s+1)\ge2^s\mathsf S_{\rm m}(s),\qquad
  2^{s+1}\mathsf S_\varrho(s+1)\ge2^s\mathsf S_\varrho(s)\qquad(0\le s<\Delta)
 \end{equation}
 holds, then $\mathsf f_{\rm m}$ and $\mathsf f_\varrho$ are nondecreasing on
 $\{0,\dots,\Delta\}$, and the bound holds with each $s_v$ replaced by any $s'_v$ with
 $s_v\le s'_v\le\Delta$.
\item Let $\mathfrak P$ be an exact partition of the nonstar subsets of $E_\times$
 into classes with fixed records whose witness digraphs have strongly connected
 components of at most two vertices, the witness sizes lying in a finite set
 $\mathcal K_{\rm w}$; for instance a refinement of the compressed partition
 (\cref{lem:advance:compressed-partition}). Let \textup{(S1)} hold at a real
 $Q>1$, and let $C_k(Q)$ be the sum over the classes of $\mathfrak P$ with witness
 size $k$ of
 $2^{\mathrm{ex}^\star(\mathcal C)/2}R_{\rm full}^2(1-\hat e)^{-k}\prod_u\mathsf
 f_{\rm m}(s_u)^{1/2}\prod_b\mathsf f_\varrho(s_b)^{1/2}$, evaluated at $r=Q^{-1/2}$.
 Then for every prime $p\ge Q$
 \[
  \norm{T_{\rm ns}}_{L^2(\mu)}\le\sum_{k\in\mathcal K_{\rm w}}C_k(Q)\,p^{-k/2}.
 \]
 Under the condition \eqref{eq:monotone-signature} the same holds when $C_k(Q)$ is
 computed from stored numbers $s'_v$ with $s_v\le s'_v\le\Delta$.
\end{enumerate}
\end{corollary}

\begin{proof}
(a) Since $t_v\le d^0_v\le\Delta^+$ and $e_t(p,n)$ is nondecreasing in $t$,
$1-e_{t_v}(p,n)\ge1-\hat e$. Let $u\in U_x$. As $H_u\subseteq C_z$, the special sum
of $F_u$ is
\[
 \sum_{b\in C_z}2^{c^U(u,b)-[b\in H_u]}
 =\sum_{b\in C_z}w^{\rm m}_{ub}-\frac12\sum_{b\in H_u}w^{\rm m}_{ub}
 \le\sum_{b\in U^+}w^{\rm m}_{ub}-\frac12\sum_{b\in H_u}w^{\rm m}_{ub}.
\]
The set $H_u\subseteq N(u)$ has $h_u=s_u-t_u\ge s_u-d^0_u$ elements, hence at least
$h^0_u(s_u)$, so its weights add up to at least $\sigma^{\rm m}_u(s_u)$; and
$d_u\ge s_u$. So the special sum is at most $\mathsf S_{\rm m}(s_u)$, and
$F_u\le\mathsf f_{\rm m}(s_u)$. For $u\in U_z$ the same argument applies with $C_x$,
and for a cover vertex $b$ with the weights $w^\varrho_{bj}$ and $H_b\subseteq N(b)$,
$|H_b|=s_b-t_b\ge h^0_b(s_b)$; this gives $F_b\le\mathsf f_\varrho(s_b)$. Finally
$2^\beta=2^{2\beta/2}$. Under this condition the numerators of $\mathsf f_{\rm m}(s)$ and
$\mathsf f_\varrho(s)$ are nondecreasing in $s$, since $e_s(p,n)$ and $e_s(p,2n)$ are
(\ref{it:et:poly}) and $r^22^s\mathsf S(s)$ is by \eqref{eq:monotone-signature}, while
the denominator does not depend on $s$; so $\mathsf f(s_v)\le\mathsf f(s'_v)$.

(b) (S1) at $Q$ implies (S1) at every prime $p\ge Q$ (\cref{def:spec:data}). By the
triangle inequality over $\mathfrak P$ and (a), $\norm{T_{\rm ns}}\le\sum_{\mathcal
C\in\mathfrak P}r^{k_{\mathcal C}}\varphi_{\mathcal C}(r)$, where
$\varphi_{\mathcal C}(r)$ is the coefficient of $r^{k_{\mathcal C}}$ in (a) as a
function of $r$. The count errors $e_s(p,n)$, $e_s(p,2n)$ and $r^2$ are polynomials
in $r$ with nonnegative coefficients, and on $(0,Q^{-1/2}]$ every $e_t(p,n')$ with
$t\le\Delta$, $n'\le2n$ is below one; so $\varphi_{\mathcal C}$ is admissible by
\cref{lem:majorant}(i),(ii) (quotients $1/(1-\varphi)$ and
$(1+\varphi)/(1-\varphi)$, products and square roots). Since the records are fixed,
$\varphi_{\mathcal C}(p^{-1/2})\le\varphi_{\mathcal C}(Q^{-1/2})$ by
\cref{lem:majorant}(a), and summing over the classes of each witness size gives
the claim. With stored numbers $s'_v$ every class bound of (a) only grows, and the
same argument applies.
\end{proof}

For a fixed partition with fixed records, the coefficients $C_k(P)$ between
two evaluation points $P_0<P_1$ are bounded by log-convexity in $\log P$. The leaf factors have a common factor;
after it is divided out, the rest is a positive combination of square roots of
positive sums of exponentials in $\log p$.

\begin{lemma}[Log-convex nonstar coefficients]\label{lem:logconvex-coefficients}
Fix $\mathfrak P$ and its records as in \cref{cor:extraction-leaf}(b), and one error
pair of the setting. For real $P>1$ at which \textup{(S1)} holds let $C_k(P)$ be the
coefficients of \cref{cor:extraction-leaf}(b), and put
$M(P)=R_{\rm full}(P)^2(1-\hat e(P))^{-n}$, both evaluated at $r=P^{-1/2}$. Then
$C_k=M D_k$, where $D_k$ is a finite positive sum of products of square roots of the
numerators of $\mathsf f_{\rm m}$ and $\mathsf f_\varrho$, and $M$, $D_k$ and $C_k$ are
nonincreasing functions of $P$ and log-convex functions of $\log P$. Consequently, if
$P_0<P_1$, \textup{(S1)} holds at $P_0$, $c_i\ge C_k(P_i)$ for $i=0,1$,
$P_0\le Q\le P_1$ and $\theta=\log(Q/P_0)/\log(P_1/P_0)$, then
\[
 C_k(Q)\le M(Q)\Bigl(\frac{c_0}{M(P_0)}\Bigr)^{1-\theta}
 \Bigl(\frac{c_1}{M(P_1)}\Bigr)^{\theta}\le c_0^{1-\theta}c_1^{\theta}.
\]
\end{lemma}

\begin{proof}
A class of witness size $k$ contributes to $C_k$ the term
$2^{\mathrm{ex}^\star/2}R_{\rm full}^2(1-\hat e)^{-k}\prod_v\mathsf f(s_v)^{1/2}$ over its
$2(n-k)$ non-witness vertices, and every $\mathsf f(s_v)$ is a numerator divided by
$1-\hat e$. So the powers of $1-\hat e$ add up to $-k-(n-k)=-n$, and the term is $M$
times $2^{\mathrm{ex}^\star/2}$ times the product of the square roots of the
numerators; this gives $C_k=MD_k$.

Put $u=\log P$, so that $r=e^{-u/2}$. Every count error of \eqref{eq:graphs:et} or
\eqref{eq:graphs:graded}, and every $r^22^s\mathsf S(s)$, is a nonnegative combination of
$e^{-u/2}$ and $e^{-u}$. A positive function of $u$ is log-convex if its logarithm is
convex; the functions $e^{cu}$ are log-convex, and products, positive powers and
positive sums of log-convex functions are log-convex, the last by H\"older's
inequality. So every numerator $1+e_s+r^22^s\mathsf S(s)$, every factor $1+e_{d^0_i}$
and $D_k$ are log-convex. If $\varepsilon$ is a nonnegative combination of decreasing
exponentials of $u$ with $\varepsilon<1$, then $(1-\varepsilon)^{-1}$ is log-convex,
since
\[
 \frac{d^2}{du^2}\bigl(-\log(1-\varepsilon)\bigr)
 =\frac{\varepsilon''}{1-\varepsilon}+\frac{(\varepsilon')^2}{(1-\varepsilon)^2}\ge0,
\]
as $\varepsilon''\ge0$. Hence $M$ and $C_k=MD_k$ are log-convex. Every exponential
$e^{-u/2}$, $e^{-u}$ decreases in $u$, so all these functions are nonincreasing in
$P$. On $[P_0,P_1]$ all denominators are positive, as (S1) at $P_0$ gives (S1) at every
$P\ge P_0$ (\cref{def:spec:data}).

Write $\log Q=(1-\theta)\log P_0+\theta\log P_1$. Log-convexity of $D_k$ gives
$D_k(Q)\le D_k(P_0)^{1-\theta}D_k(P_1)^\theta$, and $D_k(P_i)=C_k(P_i)/M(P_i)\le
c_i/M(P_i)$; multiplying by $M(Q)$ gives the first inequality. The second is the
log-convexity of $M$, $M(Q)\le M(P_0)^{1-\theta}M(P_1)^\theta$.
\end{proof}

At an evaluation point the certificates bound $C_k(P_i)$ by rounding every square
root of a vertex factor upward to a multiple of $2^{-48}$ and summing the stored
histogram exactly. The lemma is applied to the unrounded functions $C_k$, with these
upper bounds as $c_i$, so the rounded values need not be log-convex.

\begin{lemma}[An order with a bounded weighted extension sum]\label{lem:cover:order}
For an integer $g\ge0$ put $\mathrm{av}(g)=(2^{g+1}-1)/(g+1)$. If a vertex has $g$
neighbours in a uniformly ordered vertex set containing it, then the expected value
of $2^t$, where $t$ counts its earlier neighbours, is $\mathrm{av}(g)$. Consequently,
for any fixed nonnegative weights $\omega_v$ on the vertices, some deterministic
order has $\sum_v\omega_v2^{t_v}\le\sum_v\omega_v\mathrm{av}(g_v)$, where $g_v$ is
the number of neighbours of $v$ in the set.
\end{lemma}

\begin{proof}
The vertex has uniform rank among itself and its $g$ neighbours, so $t$ is uniform
on $\{0,\dots,g\}$, and $\E2^t=\frac1{g+1}\sum_{t=0}^g2^t=\mathrm{av}(g)$. By
linearity the weighted sum has expectation $\sum_v\omega_v\mathrm{av}(g_v)$, and a
finite average is at least its minimum.
\end{proof}

\begin{corollary}[Compressed classes, log-linear form]\label{cor:extraction-compressed}
Assume \textup{(S1)} at $p$. Let $\mathcal C$ be a class with a witness of size two
and a cover, and let $c_{\rm ker}(\mathcal C)\in\{1,2\}$ be its kernel constant
($2$ for the complete rectangle). For a non-witness vertex $i$ let $h_i$ be the
number of free pairs at $i$, and put
\[
\begin{aligned}
 \mathrm{oc}(\mathcal C)={}&\sum_{i\in C_x}2^{h_i}\,\mathrm{av}(|\operatorname{Adj}(i)\cap C_x|)
 +\sum_{i\in U_x}2^{h_i+d^0_i-|\operatorname{Adj}(i)\cap U_x|}\,
 \mathrm{av}(|\operatorname{Adj}(i)\cap U_x|)\\
 &+\sum_{i\in C_z}2^{h_i}\,\mathrm{av}(|\operatorname{Adj}(i)\cap C_z|)\\
 &+\sum_{i\in U_z}2^{h_i+d^0_i-|\operatorname{Adj}(i)\cap U_z|}\,
 \mathrm{av}(|\operatorname{Adj}(i)\cap U_z|),
\end{aligned}
\]
$\gamma_{\rm cmp}=nr^2+\frac\Delta4r$,
$J_{\rm sp}=\sum_{i,j\in U^+}2^{d_i+|\operatorname{Adj}(i)\cap N(j)|
+|\operatorname{Adj}(j)\cap N(i)|}$ and
\[
 \Lambda_0=\frac{R_{\rm full}^2\exp(r^2J_{\rm sp})}
 {(1-\hat e)^2\prod_{i\in U^+}(1-e_{d^0_i}(p,n))}.
\]
Then
$\norm{T_{\mathcal C}}\le c_{\rm ker}(\mathcal C)\,r^2\,2^{\mathrm{ex}(\mathcal C)/2}
\Lambda_0\exp(\gamma_{\rm cmp}\,\mathrm{oc}(\mathcal C))$.
Consequently, group the classes into \emph{buckets} indexed by
integers $\mathfrak b$: if every class of an exact partition of
the nonstar subsets is assigned a bucket index
$\mathfrak b\ge\mathrm{oc}(\mathcal C)2^{-\mathfrak q}$, for a
fixed integer
$\mathfrak q\ge0$, and
$\mathsf a_{\mathfrak b}+\mathsf b_{\mathfrak b}\sqrt2=\sum_{\mathcal C\mapsto\mathfrak b}
c_{\rm ker}(\mathcal C)2^{\mathrm{ex}(\mathcal C)/2}$ with
$\mathsf a_{\mathfrak b},\mathsf b_{\mathfrak b}\in\Z_{\ge0}$, then
\[
 \norm{T_{\rm ns}}\le r^2\Lambda_0\sum_{\mathfrak b}
 \bigl(\mathsf a_{\mathfrak b}+\mathsf b_{\mathfrak b}\sqrt2\bigr)
 \exp\bigl(\gamma_{\rm cmp}\,\mathfrak b\,2^{\mathfrak q}\bigr).
\]
For fixed classes and a fixed assignment of each class to an integer
$\mathfrak b$ (its bucket) as above, this right side divided by $r^2$ is
admissible
(\cref{lem:majorant}) as a function of $r$ on the range where \textup{(S1)} holds.
\end{corollary}

\begin{proof}
We choose the orders of the record and apply \cref{lem:matching:extractionk} with
$k=2$. For $i\in C_x$ the number $t_i$ counts the neighbours of $i$ earlier in
$C_x$; for $i\in U_x$ it counts the $d^0_i-|\operatorname{Adj}(i)\cap U_x|$
neighbours outside $U_x$ and the neighbours earlier in $U_x$. By
\cref{lem:cover:order}, applied to each of the four sets with the weights
$2^{h_i}$ (multiplied by $2^{d^0_i-|\operatorname{Adj}(i)\cap U_x|}$ on $U_x$, and
likewise on $U_z$), there are orders with $\sum_v2^{s_v}\le\mathrm{oc}(\mathcal C)$.
Since $s\le\Delta$, \ref{it:et:mono} gives
$e_s(p,n)\le e_s(p,2n)\le2^s(\frac\Delta2r+2nr^2)=2^s\cdot2\gamma_{\rm cmp}$. By the
symmetry of $N$, $\mathrm{mg}(j,i)=|\operatorname{Adj}(i)\cap N(j)|$ and
$c^U(j,i)=|\operatorname{Adj}(j)\cap N(i)|$, so $c^U(u,b)=\mathrm{mg}(b,u)$. Put
$\mathrm{sp}_v=\sum_{j\in U^+}2^{d_v+\mathrm{mg}(j,v)+c^U(j,v)}$. As $s_v\le d_v$, the
special sum of an integrated vertex $u$ is at most
$r^2\sum_b2^{s_u+\mathrm{mg}(b,u)}\le r^2\mathrm{sp}_u$, and that of a cover vertex
is at most $r^2\mathrm{sp}_b$; with $t_v\le d^0_v$, every vertex factor satisfies
\[
 F_v\le\frac{1+2\gamma_{\rm cmp}2^{s_v}+r^2\mathrm{sp}_v}{1-e_{d^0_v}(p,n)}
 \le\frac{\exp(2\gamma_{\rm cmp}2^{s_v}+r^2\mathrm{sp}_v)}{1-e_{d^0_v}(p,n)}.
\]
Each $i\in U^+$ occurs at most twice among the $2(n-2)$ non-witness vertices (once
in $X\setminus V$, once in $Z\setminus W$), and $J_{\rm sp}=\sum_i\mathrm{sp}_i$.
Hence $\prod_vF_v^{1/2}\le\exp(\gamma_{\rm cmp}\,\mathrm{oc}(\mathcal C)+r^2J_{\rm sp})
\prod_{i\in U^+}(1-e_{d^0_i}(p,n))^{-1}$, every factor $(1-e)^{-1}$ being at least
one, and \cref{lem:matching:extractionk} gives the class bound. The sum follows from
the triangle inequality, since $\mathrm{oc}(\mathcal C)\le\mathfrak b2^{\mathfrak q}$.
Finally $\gamma_{\rm cmp}$ and $r^2J_{\rm sp}$ are polynomials in $r$ with nonnegative
coefficients, and $\Lambda_0$ is built from them and from the count errors, which are
such polynomials and are below one where (S1) holds, by products, exponentials and the
quotients $1/(1-\varphi)$ and $(1+\varphi)/(1-\varphi)$; so \cref{lem:majorant}(i),(ii)
give the admissibility.
\end{proof}

\subsection{Ordered coincidences}\label{subsec:extraction:coincidences}

The collision part $T_{\rm coll}$ of \eqref{eq:record:root-kernel} lives on the
pairs of half labellings that share a label. Such a pair has an \emph{exact
coincidence matching}: the set $I$ of the $x'$-vertices whose labels are labels of
$x$, and the injective map $\phi\colon I\to U^+$ with $x'_i=x_{\phi(i)}$. For fixed
$x$ and $(I,\phi)$ the set of $x'$ with exact coincidence matching
$(I,\phi)$, called a \emph{stratum} below, is counted by labelling the
coincident vertices first;
a coincident vertex costs a factor $r^2=1/p$, and the remaining vertices are
counted sequentially in an order chosen for the stratum type.

\begin{lemma}[Ordered coincidence bound]\label{lem:oct24-collision}
Let spectral data satisfy $e_\Delta(p,2n)<1$. Put
\[
\begin{gathered}
 L_{\rm int}=\prod_{i\in U^+}\bigl(1-e_{d^0_i}(p,n)\bigr),\qquad
 \Upsilon=\prod_{i\in U^+}\bigl(1+e_{d_i}(p,2n)\bigr),\\
 a_{ij}=2^{r_i-[j\in N(i)]+\mathrm{mg}(j,i)}\quad(i,j\in U^+),\qquad
 \Xi=2^{\Delta^+/2}p^{-1}\sum_{i,j\in U^+}a_{ij},
\end{gathered}
\]
and write $T_{\rm coll}=T_{\rm coll,1}+T_{\rm coll,\ge2}$ for the parts of
$T_{\rm coll}$ on the pairs with exactly one and with at least two common labels.
\begin{enumerate}[label=\textup{(\alph*)}]
\item \emph{Exactly one coincidence.} For each $(i,j)\in U^+\times U^+$ choose an
 order of $U^+\setminus\{i\}$, and for $\ell\ne i$ let
 \[
  t^{ij}_\ell=r_\ell+[\ell\sim i]-[\ell\sim i\text{ and }j\in N(\ell)]
  +|\operatorname{Adj}(\ell)\cap\{\text{vertices before }\ell\}|.
 \]
 Then
 $\norm{T_{\rm coll,1}}\le D_1:=\dfrac1{pL_{\rm int}}\sum_{i,j\in U^+}a_{ij}
 \prod_{\ell\ne i}\bigl(1+e_{t^{ij}_\ell}(p,2n)\bigr)$.
\item \emph{At least two coincidences.} For each $k\ge2$ choose an order $\pi_k$
 of $U^+$, independently of the coincidence matching; let $b_i(\pi_k)$ be the
 number of $G^+$-neighbours of $i$ preceding $i$ in $\pi_k$,
 $t^{(k)}_i=\min\bigl(r_i+d^0_i,\ r_i+b_i(\pi_k)+k\bigr)$ and
 $\Upsilon_k=\prod_i(1+e_{t^{(k)}_i}(p,2n))$. Then
 \begin{equation}\label{eq:oct24-collision}
  \norm{T_{\rm coll,\ge2}}\le D_{\ge2}:=L_{\rm int}^{-1}\sum_{k=2}^n\Upsilon_k\,
  \frac{\Xi^k}{k!}.
 \end{equation}
\item \emph{Series tail in $k$.} $\Upsilon_k=\Upsilon$ for
 $k\ge\Delta^+$. Put
 $k_0=\max(2,\Delta^+)$. If $\Xi<k_0+1$, then
 \[
  \frac1{L_{\rm int}}\sum_{k\ge k_0}\Upsilon_k\frac{\Xi^k}{k!}\le
  \frac\Upsilon{L_{\rm int}}\,\frac{\Xi^{k_0}}{k_0!}\Bigl(1-\frac\Xi{k_0+1}\Bigr)^{-1}.
 \]
\item \emph{Coarse forms.} For all orders,
 $D_1\le\Upsilon L_{\rm int}^{-1}p^{-1}\sum_{i,j}a_{ij}$ and
 $D_{\ge2}\le\Upsilon L_{\rm int}^{-1}(e^\Xi-1-\Xi)$. In particular
 \begin{equation}\label{eq:collision-natural}
  \norm{T_{\rm coll}}\le D_{\rm nat}:=\frac\Upsilon{L_{\rm int}}
  \Bigl[p^{-1}\sum_{i,j\in U^+}a_{ij}+e^\Xi-1-\Xi\Bigr],
 \end{equation}
 the \emph{natural bound}.
\item \emph{Minima and endpoint monotonicity.} Put
 $D_{\rm ord}=D_1+\Upsilon L_{\rm int}^{-1}(e^\Xi-1-\Xi)$, the ordered
 single-coincidence bound. The minimum of any finitely many valid bounds for
 $\norm{T_{\rm coll}}$ --- in particular of $D_1+D_{\ge2}$ (with the tail of
 \textup{(c)} for $k\ge k_0$), of $D_{\rm ord}$ and of the natural bound
 $D_{\rm nat}$, for any choices of the orders --- is again a bound for
 $\norm{T_{\rm coll}}$. If the orders are fixed and $e_\Delta(P,2n)<1$ at a real
 $P>1$, then at every prime $p\ge P$ the quantities $D_1/r^2$, $D_{\ge2}/r^4$,
 $D_{\rm ord}/r^2$ and $D_{\rm nat}/r^2$ are at most their values at $P$, and so is
 the tail of \textup{(c)} divided by $r^4$ if $\Xi<k_0+1$ at $P$.
\end{enumerate}
\end{lemma}

\begin{proof}
The kernel $K_{\rm coll}$ is nonnegative and symmetric (\cref{lem:spec:expansion}),
and the number of common labels is symmetric too; so both parts have nonnegative
symmetric kernels, and by Schur's test their norms are
at most the supremum over $x$ of their row sums $\E_{x'\sim\mu}[\cdot]$. Fix
$x\in\mathcal I^+$. By \cref{lem:graphs:sequential}\ref{it:seq:count} with the error
pair of the setting and any order, $|\mathcal I^+|\ge p^n2^{-|E(G^+)|}L_{\rm int}$, using
$1-e_{t}(p,n)\ge1-e_{d^0_i}(p,n)$ for $t\le d^0_i$. In a stratum $(I,\phi)$ the
cross factors at an $x'$-vertex $i$ are $w_i(x'_i)$ with
$w_i(y)=\prod_{j\in N(i)}(1+\chi(\varsigma_{ji}(y-x_j)))$; at a coincident vertex,
$x'_i=x_{\phi(i)}$, the factor of $j=\phi(i)$, if $\phi(i)\in N(i)$, is $1+\chi(0)=1$
and the others are at most $2$, so $w_i(x_{\phi(i)})\le2^{r_i-[\phi(i)\in N(i)]}$.
We drop the internal constraints among coincident vertices, which only increases
the sums. The stratum is summed by labelling the vertices of $x'$ one at a time;
the sum of $w_i$ over the admissible labels of $i$, given the earlier labels, is
the \emph{one-step sum} at $i$, and the stratum sum is at most the product of
bounds for the one-step sums that are uniform in the earlier labels (sum the
last vertex first).

(a) Let $I=\{i\}$ and $\phi(i)=j$. Label $i$ first and the other vertices in the
chosen order. A vertex $\ell\ne i$ has label outside $x(U^+)$, so its cross roots
$x_{j'}$ ($j'\in N(\ell)$) and its internal roots (the labels of its labelled
neighbours: $x_j$ if $\ell\sim i$, and labels outside $x(U^+)$ otherwise) merge
exactly when $\ell\sim i$ and $j\in N(\ell)$. With $\mathcal X=x(U^+)\cup\{$labels
of $x'$ already chosen$\}$, \cref{lem:extraction-merged-count} bounds its one-step sum by
$p2^{-t^0_\ell}2^{\mu_\ell}(1+e_{t^{ij}_\ell}(p,2n))$, where $t^0_\ell$ is its number
of labelled neighbours, $\mu_\ell=[\ell\sim i\text{ and }j\in N(\ell)]$, and
$t^{ij}_\ell=r_\ell+t^0_\ell-\mu_\ell$ is its number of distinct roots. Every edge of
$G^+$ is counted in $t^0_\ell$ at its later endpoint, so $\sum_{\ell\ne
i}t^0_\ell=|E(G^+)|$, and $\sum_{\ell\ne i}\mu_\ell=\mathrm{mg}(j,i)$. The stratum sum
is therefore at most
$p^{n-1}2^{-|E(G^+)|}a_{ij}\prod_{\ell\ne i}(1+e_{t^{ij}_\ell}(p,2n))$, and dividing by
$|\mathcal I^+|$ and summing over $(i,j)$ bounds the row sum by $D_1$.

(b) Fix a stratum $(I,\phi)$ with $|I|=k\ge2$, label $I$ first and the other vertices
in the restriction of $\pi_k$. A vertex $i\notin I$ has $r_i$ distinct cross roots
and $t^0_i$ internal roots, the labels of its labelled neighbours, where
$t^0_i\le\min(d^0_i,b_i(\pi_k)+k)$: the labelled neighbours are at most $k$
coincident ones and the neighbours preceding $i$ in $\pi_k$. A cross root $x_j$ equals
an internal root exactly when a coincident neighbour $l\in I$ has $\phi(l)=j\in N(i)$;
let $\mu_i$ be the number of such $l$. As $r_i+t^0_i-\mu_i\le t^{(k)}_i$,
\cref{lem:extraction-merged-count} bounds the one-step sum at $i$ by
$p2^{-t^0_i}2^{\mu_i}(1+e_{t^{(k)}_i}(p,2n))$, whatever the earlier labels. Each
merge counted in $\mu_i$ is caused by some $l\in I$ with $i\in\operatorname{Adj}(l)$
and $\phi(l)\in N(i)$, so
$\sum_{i\notin I}\mu_i\le\sum_{l\in I}\mathrm{mg}(\phi(l),l)$. Every edge of $G^+$
not inside $I$ is counted in $t^0$ of its later endpoint, which lies outside $I$;
so $\sum_{i\notin I}t^0_i=|E(G^+)|-e(I)$. Inserting the factors of $\Upsilon_k$ at
the coincident vertices, each at least one, the stratum sum is at most
$p^{n-k}2^{-|E(G^+)|+e(I)}\prod_{l\in I}a_{l,\phi(l)}\Upsilon_k$, and after division by
$|\mathcal I^+|$ at most $L_{\rm int}^{-1}\Upsilon_kp^{-k}2^{e(I)}\prod_{l\in
I}a_{l,\phi(l)}$. Now $2^{e(I)}\le2^{k\Delta^+/2}$, and expanding
$(\sum_{i,j}a_{ij})^k$ counts every pair $(I,\phi)$ with $|I|=k$ exactly $k!$ times,
besides other nonnegative terms. Summing over the strata with $|I|=k$ and over
$k\ge2$ bounds the row sum by $D_{\ge2}$.

(c) For $k\ge\Delta^+$, $b_i(\pi_k)+k\ge d^0_i$, so $t^{(k)}_i=r_i+d^0_i=d_i$. The
ratio of consecutive terms of $\sum_{k\ge k_0}\Xi^k/k!$ is at most $\Xi/(k_0+1)<1$,
which gives the geometric bound.

(d) Always $t^{ij}_\ell\le r_\ell+d^0_\ell=d_\ell$ and $t^{(k)}_i\le d_i$, so every
product is at most $\Upsilon$, and $\sum_{k\ge2}\Xi^k/k!\le e^\Xi-1-\Xi$.

(e) Each listed quantity bounds the same norm, so their minimum does; no error
bound is subtracted. For the endpoints write $r=p^{-1/2}$. Every count error is a
polynomial in $r$ with nonnegative coefficients, $L_{\rm int}^{-1}$ is a product of
factors $1/(1-\varphi)$ with $\varphi<1$ on $(0,P^{-1/2}]$, and $\Xi=r^2\Xi_0$ with a
constant $\Xi_0$. Hence $D_1/r^2$, $D_{\ge2}/r^4=L_{\rm int}^{-1}\sum_k\Upsilon_k\Xi_0^k
r^{2k-4}/k!$, $(e^\Xi-1-\Xi)/r^2=\sum_{k\ge2}\Xi_0^kr^{2k-2}/k!$, and so
$D_{\rm ord}/r^2$ and $D_{\rm nat}/r^2$, and the tail divided by $r^4$ are admissible
in the sense of \cref{lem:majorant}, and \cref{lem:majorant}(a) gives the claim.
\end{proof}

\begin{remark}[Why the coincident vertices are labelled first]\label{rem:collision-fixed-order}
With one fixed order of all vertices for every stratum, a coincident vertex $i$
pays $2^{t^0_i}\le2^{d^0_i}$ for its earlier neighbours, that is, the charge
$r^22^{d_i-[j\in N(i)]+\mathrm{mg}(j,i)}=2^{d^0_i}r^2a_{ij}$ instead of $r^2a_{ij}$.
Labelling the coincident vertices first moves these factors into the single factor
$2^{e(I)}$, which is $1$ for one coincidence. This factor cannot be dropped: the
internal constraint between two coincident vertices is not counted, although the
normalization $|\mathcal I^+|$ contains a factor $\frac12$ for it. The
cardinality-sensitive orders of (b) refine the count of the vertices outside $I$
once more: a vertex has at most $k$ coincident neighbours, so in a stratum with
few coincidences its internal degree is effectively that of the order $\pi_k$.
\end{remark}

\paragraph{Soft factors and incident weights.}
The dense comparisons of the upper range (\cref{prop:dense-comparison}) use
kernels whose cross factors are
$1+D_e\chi_e$ with $|D_e|\le1$, and the L20 cut of the middle root (\cref{prop:l20-collisions}) needs the
coincidence bound with the actual cross factors of the coincident vertices. Both
follow from the proofs above.

\begin{corollary}[Soft cross factors]\label{cor:soft-kernels}
Let numbers $D_e\in[-1,1]$ be given for $e\in E_\times$, put
$D_S=\prod_{e\in S}D_e$, and consider the \emph{soft kernel}
$K^D_0(x,x')=\one\{x(U^+)\cap x'(U^+)=\emptyset\}\prod_{e\in E_\times}(1+D_e\chi_e(x,x'))$,
its nonstar part $T^D_{\rm ns}=\sum_{S\ \rm nonstar}D_ST_S$, and its collision part
$T^D_{\rm coll}$, the operator with kernel
$\one\{x(U^+)\cap x'(U^+)\ne\emptyset\}\prod_{e\in E_\times}(1+D_e\chi_e)$; as in
\cref{lem:spec:expansion},
$K^D_0=\one\otimes\one+\sum_{S\ \rm star}D_ST_S+T^D_{\rm ns}-T^D_{\rm coll}$ on
$\mathcal I^+\times\mathcal I^+$.
\begin{enumerate}[label=\textup{(\alph*)}]
\item For a class $\mathcal C$ with a record, $T^D_{\mathcal C}=\sum_{S\in\mathcal C}D_ST_S$
 satisfies the bounds of \cref{lem:matching:extractionk}, \cref{cor:extraction-leaf}(a)
 and \cref{cor:extraction-compressed}; hence $T^D_{\rm ns}$ satisfies those of
 \cref{cor:extraction-leaf}(b) and \cref{cor:extraction-compressed}, with the
 partitions, records and constants computed for $E_\times$.
\item $T^D_{\rm coll}$, and its parts on the pairs with exactly one and with at least
 two common labels, satisfy the bounds of \cref{lem:oct24-collision}.
\item The same holds for the pinned rows of \cref{subsec:push:pinnedremainder}: the
 bounds of \cref{prop:pinned-extraction,prop:pinned-collision} hold when the cross
 factors are $1+D_{ij}\chi_{ij}$ with $|D_{ij}|\le1$.
\end{enumerate}
In particular, let $E'\supseteq\{e\in E_\times:D_e\ne0\}$ be the set of cross pairs of
spectral data with the same half graph $G^+$. As $D_e\chi(\varsigma_ey)=D'_e\chi(\varsigma'_ey)$
with $D'_e=\pm D_e$ for any two signs $\varsigma_e,\varsigma'_e$, the soft kernel is a soft
kernel of $E'$ (with the signs of $E'$ and $D'_e=0$ off the support of $D$). So the
bounds computed for the Boolean kernel of $E'$, all of whose coefficients are one,
bound its nonstar and collision parts, and bounds certified at a threshold $p_0$ hold
at every prime $p\ge p_0$ (\cref{cor:extraction-leaf}(b), \cref{lem:oct24-collision}(e)).
\end{corollary}

\begin{proof}
The expansion is $\prod_e(1+D_e\chi_e)=\sum_SD_S\chi_S$, and the identity follows as in
\cref{lem:spec:expansion}. The proofs of the cited statements use the cross factors
only in four ways, and each survives the soft factors. (1) The class sum is
$\sum_{S\in\mathcal C}D_S\chi_S=\prod_{e\in P}D_e\chi_e\prod_{e\in F}(1+D_e\chi_e)$: a
present pair contributes a factor of modulus at most one, so the witness kernel is
multiplied by a constant of modulus at most one, and a free pair contributes a factor
in $[0,2]$. (2) Every bound for a sum over one label of a product of free factors is
uniform in the prescribed signs: $1+b\chi=\frac{1+b}2(1+\chi)+\frac{1-b}2(1-\chi)$ is a
convex combination of the two orientations, so a product of such factors is a convex
combination of products of factors $1\pm\chi(\cdot)$, each of which
\cref{lem:extraction-merged-count} bounds by the same quantity, whatever the signs. No
free factor is squared: the Cauchy--Schwarz steps use the free factors as weights (Step
5 of \cref{lem:matching:extractionk}). (3) At a special value or at a coincidence the
factor of the coincident pair is $1+b\chi(0)=1$, and every other factor is at most
$1+|b|\le2$. (4) The collision kernel is nonnegative. It need not be symmetric, and
Schur's test is then applied with its row sums and its column sums; the column sums are
the row sums of the soft kernel with the coefficient $D_{(j,k)}$ at the pair $(k,j)$,
which is a soft kernel of the same spectral data (as in the proof of
\cref{lem:spec:expansion}), so the same bound holds for both. The pinned rows are
covered by the same four remarks; their collision bound already uses row and column
sums. Only absolute values enter: every term of the expansion is bounded, and no
monotonicity of a signed norm under the deletion of pairs is used.
\end{proof}

\begin{corollary}[Incident weights at coincident vertices]\label{cor:incident-weights}
In \cref{lem:oct24-collision}, with hard or soft cross factors
(\cref{cor:soft-kernels}), fix $x\in\mathcal I^+$. For an $x'$-vertex $i$ let
$w_i(y)=\prod_{l\in N(i)}\bigl(1+D_{(l,i)}\chi(\varsigma_{li}(y-x_l))\bigr)$ be its cross
factor at the label $y$ ($D\equiv1$ for the hard kernel), and put
$\alpha_{ij}(x)=w_i(x_j)2^{\mathrm{mg}(j,i)}$ and
$\rho_i(x)=r^2\sum_{j\in U^+}\alpha_{ij}(x)$. Then $\alpha_{ij}(x)\le a_{ij}$, and with the
orders of \cref{lem:oct24-collision}(a),(b) the row sums at $x$ of $T_{\rm coll,1}$
and of $T_{\rm coll,\ge2}$ are at most
\[
 \frac1{L_{\rm int}}\sum_{i,j\in U^+}r^2\alpha_{ij}(x)
 \prod_{\ell\ne i}\bigl(1+e_{t^{ij}_\ell}(p,2n)\bigr)
 \quad\text{and}\quad
 \frac1{L_{\rm int}}\sum_{k=2}^n\Upsilon_k\sum_{I\subseteq U^+,\,|I|=k}2^{e(I)}
 \prod_{i\in I}\rho_i(x).
\]
If the kernel is symmetric, the norms of $T_{\rm coll,1}$ and $T_{\rm coll,\ge2}$ are
at most the suprema over $x$ of these bounds; in general they are at most the
geometric means of these suprema and of the same suprema for the transposed
coefficients. A stratum $(I,\phi)$ with $\alpha_{i\phi(i)}(x)=0$ for some $i\in I$
contributes nothing, and any relation among the numbers $\alpha_{ij}(x)$ that holds for
every $x$ (vanishing weights, disjoint supports, a bound for the number of nonzero
ones) may be used to bound these sums before the supremum is taken.
\end{corollary}

\begin{proof}
In the proof of \cref{lem:oct24-collision} keep, at a coincident vertex $i$ with
label $x_{\phi(i)}$, its cross factor $w_i(x_{\phi(i)})$ instead of its bound
$2^{r_i-[\phi(i)\in N(i)]}$; that bound is $w_i(x_j)\le2^{r_i-[j\in N(i)]}$, the factor
of $l=j$ being $1+D_{(j,i)}\chi(0)=1$ and every other factor at most $2$, which gives
$\alpha_{ij}\le a_{ij}$. The one-step sums at the vertices outside $I$ are bounded as there
(\cref{cor:soft-kernels}). In (a) the stratum $(\{i\},j)$ then has sum at most
$p^{n-1}2^{-|E(G^+)|}\alpha_{ij}(x)\prod_{\ell\ne i}(1+e_{t^{ij}_\ell}(p,2n))$. In (b)
the stratum $(I,\phi)$ has, after division by $|\mathcal I^+|$, sum at most
$L_{\rm int}^{-1}\Upsilon_k2^{e(I)}\prod_{l\in I}r^2\alpha_{l\phi(l)}(x)$, and the sum
over the injective maps $\phi\colon I\to U^+$ is at most the sum over all maps, which
is $\prod_{l\in I}\rho_l(x)$. These bounds hold pointwise in $x$, which gives the last
assertion, and Schur's test is applied as in \cref{lem:oct24-collision}, with the
column sums of \cref{cor:soft-kernels} for a nonsymmetric kernel. The balanced
coincidence bound of the L20 cut is this corollary (\cref{prop:l20-collisions}).
\end{proof}

\begin{lemma}[Sign-incompatible collision strata]\label{lem:signed-collisions}
Suppose $p\equiv3\pmod4$ and the two half graphs have the same signed
adjacency $A_{ij}=\one[i\to j]-\one[j\to i]$. Write the hard cross
kernel as $\prod_{k,i}(1+D_{ki}\chi(y_i-x_k))$, omitting zero
coefficients, with $D_{ki}\in\{-1,0,1\}$. Define
\[
 \gamma_{ij}=\one[A_{kj}D_{ki}\ne-1\text{ for all }k,\quad
                       A_{i\ell}D_{j\ell}\ne-1\text{ for all }\ell].
\]
Every collision stratum containing $y_i=x_j$ vanishes when
$\gamma_{ij}=0$. Consequently in the nonnegative collision sums of
\cref{lem:oct24-collision}, with the notation of that lemma, one may replace the old pin weights by
\[
 \widetilde a_{ij}=\gamma_{ij}
       2^{r_i-\one[j\in N(i)]+\mathrm{mg}(j,i)},\qquad
 s_i=\sum_j\widetilde a_{ij}.
\]
In particular, for the same fixed insertion orders, the single- and
multiple-collision row bounds become
\begin{align*}
 D_1^\#&=\frac q{L_{\rm int}}\sum_{i,j}\widetilde a_{ij}
                      \prod_{\ell\ne i}(1+e_{t^{ij}_\ell}(p,2n)),\\
 D_{\ge2}^\#&=\frac1{L_{\rm int}}\sum_{k=2}^n q^k\Upsilon_k
                  \sum_{|I|=k}2^{e(I)}\prod_{i\in I}s_i.
\end{align*}
The transposed argument supplies the column bounds.
\end{lemma}
\begin{proof}
On the pin $y_i=x_j$, an internal plus constraint fixes
$\chi(x_j-x_k)=A_{kj}$. If its product with $D_{ki}$ is $-1$,
the corresponding cross factor is zero. The second test uses the
internal minus constraint
$\chi(y_\ell-y_i)=A_{i\ell}$ in exactly the same way.
Within-half injectivity guarantees that these differences are nonzero.
Adding further cross coincidences cannot remove the vanishing factor.
Delete these zero strata before applying the positive extension bounds.
For multiple collisions, sum first over injective compatible target
maps, then enlarge to all compatible maps; the latter sum factors
as $\prod_{i\in I}s_i$. The internal pinned-edge cost $2^{e(I)}$
remains. This proves both bounds.
\end{proof}

The second sign test is a property of the entire stratum; it is not
an upper bound for an isolated incident factor. Nor does it apply to
a softened factor $1+b\chi$ with $|b|<1$. These distinctions matter
when the same graph is used for different comparison laws.

\subsection{Conditional laws given fixed labels}\label{subsec:push:pinnedremainder}

The conditional rows of the upper range (\cref{subsec:advance:product-rows}) are
labellings of $n$ labels that avoid $z$ fixed labels and satisfy a few square-class
constraints against fixed labels; there are no edges inside a row, and the cross
constraints between two rows need not be symmetric. The pins play the role of the
half graph: a fixed root never coincides with a free root, so there are no merge
charges, while every count gains the pins and the fixed labels as exclusions. Before
the rows are made injective their labels are independent, each uniform on its allowed
set; so the star part of a row factors through one coefficient matrix, as the stars of
an invariant law do in \cref{cor:star-matrix} (\cref{prop:conditional-star}), and the
fold and matching rows of one coordinate, which have the same allowed sets, are
bounded jointly on one chart (\cref{lem:joint-rows}).

\paragraph{Setting.}
Let $n\ge2$, $z\ge0$, and let $\mathcal M\subset\Fp$ be a set of $z$ fixed labels.
Every coordinate $i\in[n]$ carries a set of $\varpi_i\le\varpi$ \emph{pins}
$(\beta,\epsilon)\in\mathcal M\times\{\pm1\}$ with distinct roots $\beta$, the same
for both copies. Let $\mathcal I_{\rm pin}$ be the set of injective
$x\in(\Fp\setminus\mathcal M)^n$ with $\epsilon(x_i-\beta)\in\Qsq$ for every pin
$(\beta,\epsilon)$ of every $i$, and let $\mu$ be the uniform law on
$\mathcal I_{\rm pin}$. The \emph{allowed set} of a coordinate $i$ consists of the
$t\in\Fp\setminus\mathcal M$ with $\epsilon(t-\beta)\in\Qsq$ for every pin
$(\beta,\epsilon)$ of $i$; with $\lambda_i$ the uniform law on it and
$\lambda=\bigotimes_i\lambda_i$, the law $\mu$ is $\lambda$ conditioned on
injectivity. Let
$E_\times\subseteq[n]\times[n]$ carry signs $\varsigma_{ij}\in\{\pm1\}$, and put
$N_X(i)=\{j:(i,j)\in E_\times\}$ and $N_Y(j)=\{i:(i,j)\in E_\times\}$; no symmetry
is assumed. The rows are the copies $X$ (labels $x$) and $Y$ (labels $y$) of $[n]$;
$\chi_{ij}(x,y)=\chi(\varsigma_{ij}(y_j-x_i))$, and $T_S$, classes, witnesses and
covers are as in \cref{subsec:extraction:partition}, with $y$ in place of $x'$. For a
non-witness vertex $v$ let $h_v$ be its number of free pairs, counted in $N_X$ on
the row side and in $N_Y$ on the column side. Put
\[
\begin{gathered}
 e^{\rm pin}_d=2^d\Bigl((2n+z)r^2+\frac d2\,r\Bigr),\qquad
 \Delta_{\rm pin}=\varpi+\max_{i\in[n]}\max\bigl(|N_X(i)|,|N_Y(i)|\bigr),\\
 \gamma_{\rm pin}=\frac{2n+z}2\,r^2+\frac{\Delta_{\rm pin}}4\,r,\qquad
 R_{\rm pin}=\Bigl(\frac{1+e^{\rm pin}_\varpi}{1-e^{\rm pin}_\varpi}\Bigr)^n,\qquad
 \Lambda^{\rm pin}_0=\frac{R_{\rm pin}^2\exp(r^2n^22^{\Delta_{\rm pin}})}
 {(1-e^{\rm pin}_\varpi)^{n+2}}.
\end{gathered}
\]
By \ref{it:et:mono}, $e_d(p,n')\le e^{\rm pin}_d$ for $n'\le2n+z$; so by
\cref{lem:extension}\ref{it:ext:upper},\ref{it:ext:coarse} every set $\Phi$ of
\cref{lem:extension} with $d$ roots and at most $2n+z$ excluded labels has between
$p2^{-d}(1-e^{\rm pin}_d)$ and $p2^{-d}(1+e^{\rm pin}_d)$ elements, and the upper
bound needs no condition on the exclusions.

\begin{proposition}[Pinned compressed classes]\label{prop:pinned-extraction}
Assume $e^{\rm pin}_\varpi<\frac12$. Let $\mathcal C$ be a class of nonstar subsets of
$E_\times$ with a witness of size two whose kernel, compressed to injective pairs,
has norm at most $c_{\rm ker}(\mathcal C)r^2$, $c_{\rm ker}(\mathcal C)\in\{1,2\}$,
and with a cover; put $\mathrm{ex}^{\mathrm{pin}}(\mathcal C)=f_V+f_W$ (the pinned analogues of $\mathrm{ex}$ and $\mathrm{oc}$ in \cref{cor:extraction-compressed}, with different definitions) and
$\mathrm{oc}^{\mathrm{pin}}(\mathcal C)=\sum_v2^{h_v}$ over the $2(n-2)$ non-witness vertices. Then
\begin{equation}\label{eq:push:pinnedclass}
 \norm{T_{\mathcal C}}_{L^2(\mu)\to L^2(\mu)}
 \le c_{\rm ker}(\mathcal C)\,r^2\,2^{\mathrm{ex}^{\mathrm{pin}}(\mathcal C)/2+2\varpi}\,
 \Lambda^{\rm pin}_0\exp\bigl(\gamma_{\rm pin}2^\varpi\,\mathrm{oc}^{\mathrm{pin}}(\mathcal C)\bigr).
\end{equation}
Consequently, for an exact partition of the nonstar subsets of $E_\times$ into such
classes (for instance the compressed partition of
\cref{lem:advance:compressed-partition}(b), with $c_{\rm ker}$ from
\cref{lem:matching:kernel}), grouped by integers
$\mathfrak b\ge\mathrm{oc}^{\mathrm{pin}}(\mathcal C)2^{-\mathfrak q}$ with
$\mathsf a_{\mathfrak b}+\mathsf b_{\mathfrak b}\sqrt2=\sum_{\mathcal C\mapsto\mathfrak b}
c_{\rm ker}(\mathcal C)2^{\mathrm{ex}^{\mathrm{pin}}(\mathcal C)/2}$, the nonstar part is bounded by
\[
 N^{\rm pin}=r^2\Lambda^{\rm pin}_0\,2^{2\varpi}\sum_{\mathfrak b}
 \bigl(\mathsf a_{\mathfrak b}+\mathsf b_{\mathfrak b}\sqrt2\bigr)
 \exp\bigl(\gamma_{\rm pin}2^\varpi\,\mathfrak b\,2^{\mathfrak q}\bigr).
\]
\end{proposition}

\begin{proof}
We follow the proof of \cref{lem:matching:extractionk} for $\mathcal C$, with
$k=2$, the pinned law $\mu$ on both rows, and any orders; only the counts change.
Every count below is that of a set $\Phi$ of \cref{lem:extension} with at most
$2n+z$ excluded labels.

Steps 1, 2 and 5 use only the product form of the weights, the cover condition
and the independence of the two rows under $\mu\otimes\mu$, and hold verbatim.

\emph{Step 3.} Given the cover labels $c=x_{C_x}$, a label of $U_x$ meets only its
$\varpi_u\le\varpi$ pins, whose roots lie in $\mathcal M$ and so are distinct, and
avoids $\mathcal M$ and the labels already used. So the number of completions of
$(c,v)$ on $U_x$ lies in $\prod_{u\in U_x}p2^{-\varpi_u}[1-e^{\rm pin}_\varpi,
1+e^{\rm pin}_\varpi]$, a ratio at most $R_{\rm pin}$; in particular every admissible
$(c,v)$ has a completion, as $e^{\rm pin}_\varpi<1$. Labelling the two witness
labels in turn, there are at least $p^22^{-\varpi_{j_1}-\varpi_{j_2}}
(1-e^{\rm pin}_\varpi)^2$ admissible pairs $v$. Hence $\nu_V$ is supported on
injective pairs and has density at most $D_V=2^{2\varpi}R_{\rm pin}
(1-e^{\rm pin}_\varpi)^{-2}$ with respect to uniform probability on $\Fp^2$, and
likewise $D_W$. This replaces $2^{t_V}R_{U_x}(1-\hat e)^{-2}$: each witness label
contributes the factor $2^{\varpi}$ of its pins.

\emph{Step 4} holds verbatim: the kernel bound for uniform probability, compressed to
injective pairs, applies to the zero extensions of $\mathsf a\tilde\nu_V$ and
$\mathsf b\tilde\nu_W$.

\emph{Step 6.} An integrated vertex $u\in U_x$ meets $\varpi_u$ pins and the $h_u$
cross roots $c'_b$ ($b\in H_u\subseteq C_z$). The cross roots are labels of the other
row and avoid $\mathcal M$, so no root merges, and there are no internal edges. For the
generic values $y\notin c'(C_z)$, \cref{lem:extraction-merged-count} with the pins as
the conditions $\chi(y-\beta_i)=s_i$ bounds the contribution by
$p2^{-\varpi_u}(1+e^{\rm pin}_{\varpi_u+h_u})$; the at most $n$ special values
$y=c'_b$ contribute at most $2^{h_u}=p2^{-\varpi_u}r^22^{\varpi_u+h_u}$ each. The lower
count is $|\mathcal Y|\ge\prod_up2^{-\varpi_u}(1-e^{\rm pin}_\varpi)$. So
$\mathsf m\le\prod_{u\in U_x}F^{\rm pin}_u$ with
$F^{\rm pin}_v=(1+e^{\rm pin}_{\varpi_v+h_v}+r^2n2^{\varpi_v+h_v})/(1-e^{\rm pin}_\varpi)$.

\emph{Step 7.} The cover rows are bounded as in \cref{lem:extraction-rows}, without
merge weights: the marginal of $\mu$ on $C_z$ has density at most $R_{\rm pin}$ with
respect to the uniform law on the pinned injective labellings of $C_z$, which has at
least $\prod_{b\in C_z}p2^{-\varpi_b}(1-e^{\rm pin}_\varpi)$ elements; a value
$y\notin x([n])$ of $b$ contributes as a generic value above, and the at most $n$
values $y=x_j$ at most $2^{h_b}=p2^{-\varpi_b}r^22^{\varpi_b+h_b}$ each. So the cover
row of $b$ costs $F^{\rm pin}_b$, with $h_b$ the number of free pairs at $b$ (in $N_Y$).

\emph{Step 8.} The transposed class is a class of the same kind for the support
$\{(j,i):(i,j)\in E_\times\}$, whose row and column neighbourhoods are $N_Y$ and
$N_X$, with the same pinned law on both rows. Steps 5--7 apply to it with its own
counts, which are the counts $h_v$ of the statement on $U_z$ and $C_x$. The symmetry
of the support is never used.

\emph{Conclusion.} As in Step 8 of \cref{lem:matching:extractionk},
$\norm{T_{\mathcal C}}\le c_{\rm ker}r^2\sqrt{D_VD_W}\,2^{(f_V+f_W)/2}R_{\rm pin}
\prod_vF^{\rm pin}_v{}^{1/2}$. Since $\varpi_v+h_v\le\Delta_{\rm pin}$ and
$e^{\rm pin}_d$ is nondecreasing in $d$,
$e^{\rm pin}_{\varpi_v+h_v}\le e^{\rm pin}_{\varpi+h_v}\le2^\varpi2^{h_v}\cdot2\gamma_{\rm pin}$
and $r^2n2^{\varpi_v+h_v}\le r^2n2^{\Delta_{\rm pin}}$. With $1+x\le e^x$ and the
$2(n-2)\le2n$ factors,
$\prod_vF^{\rm pin}_v{}^{1/2}\le\exp(\gamma_{\rm pin}2^\varpi\,\mathrm{oc}^{\mathrm{pin}}(\mathcal C)
+r^2n^22^{\Delta_{\rm pin}})(1-e^{\rm pin}_\varpi)^{-n}$, and
$\sqrt{D_VD_W}=2^{2\varpi}R_{\rm pin}(1-e^{\rm pin}_\varpi)^{-2}$. Collecting the
factors gives \eqref{eq:push:pinnedclass}, and the sum follows as in
\cref{cor:extraction-compressed}.
\end{proof}

\begin{proposition}[Pinned coincidences]\label{prop:pinned-collision}
In the setting of this subsection assume $e^{\rm pin}_\varpi<1$, and let
$T^{\rm pin}_{\rm coll}$ be the operator on $L^2(\mu)$ with kernel
$\one\{x([n])\cap y([n])\ne\emptyset\}\prod_{(i,j)\in E_\times}(1+\chi_{ij}(x,y))$
relative to $\mu\otimes\mu$. Put
\[
\begin{gathered}
 U_Y=r^2\sum_{j,i\in[n]}2^{\varpi+|N_Y(j)|-[i\in N_Y(j)]},\qquad
 \Upsilon_Y=\prod_{j\in[n]}\bigl(1+e^{\rm pin}_{\varpi+|N_Y(j)|}\bigr),\\
 D_Y=\frac{\Upsilon_Y}{(1-e^{\rm pin}_\varpi)^n}\bigl(e^{U_Y}-1\bigr),
\end{gathered}
\]
and let $U_X$, $\Upsilon_X$, $D_X$ be the same with $N_X$ in place of $N_Y$. Then
$\norm{T^{\rm pin}_{\rm coll}}\le(D_XD_Y)^{1/2}\le\max(D_X,D_Y)$. If the multisets
$\{|N_X(i)|\}_i$ and $\{|N_Y(j)|\}_j$ coincide, then $D_X=D_Y$. For fixed data,
$D_X/r^2$, $D_Y/r^2$ and $N^{\rm pin}/r^2$ of \cref{prop:pinned-extraction} are
admissible (\cref{lem:majorant}) as functions of $r$ on the range where
$e^{\rm pin}_\varpi<1$.
\end{proposition}

\begin{proof}
The kernel is nonnegative, so by Schur's test its norm is at most the geometric mean
of the suprema of its row sums $\E_{y\sim\mu}[\cdot]$ and its column sums
$\E_{x\sim\mu}[\cdot]$. Fix $x\in\mathcal I_{\rm pin}$ and a coincidence pattern: a
nonempty set $I$ of $y$-vertices and a map $\phi\colon I\to[n]$ with $y_j=x_{\phi(j)}$
for $j\in I$ and $y_j\notin x([n])$ otherwise. The cross factors at $j$ are
$\prod_{i\in N_Y(j)}(1+\chi(\varsigma_{ij}(y_j-x_i)))$. A coincident vertex $j$ has
weight at most $2^{|N_Y(j)|-[\phi(j)\in N_Y(j)]}=p2^{-\varpi_j}\,r^2
2^{\varpi_j+|N_Y(j)|-[\phi(j)\in N_Y(j)]}$. A free vertex $j$ meets its pins and the
$|N_Y(j)|$ cross roots $x_i$, which avoid $\mathcal M$; with no internal edges no root
merges, and \cref{lem:extraction-merged-count} bounds its weight by
$p2^{-\varpi_j}(1+e^{\rm pin}_{\varpi+|N_Y(j)|})$, whatever the order. Dividing by
$|\mathcal I_{\rm pin}|\ge\prod_jp2^{-\varpi_j}(1-e^{\rm pin}_\varpi)$, the stratum
contributes at most
$(1-e^{\rm pin}_\varpi)^{-n}\prod_{j\in I}u_{j,\phi(j)}\prod_{j\notin I}
(1+e^{\rm pin}_{\varpi+|N_Y(j)|})$, where $u_{j,i}=r^22^{\varpi+|N_Y(j)|-[i\in N_Y(j)]}$.
Summing over all nonempty $I$ and all maps $\phi$ (an upper bound), and with
$u_j=\sum_iu_{j,i}$, the row sum is at most
$\Upsilon_Y(1-e^{\rm pin}_\varpi)^{-n}(\prod_j(1+u_j)-1)$, and
$\prod_j(1+u_j)-1\le e^{U_Y}-1$. So every row sum is at most $D_Y$.
Exchanging the two rows exchanges $N_X$ and $N_Y$ and leaves the pinned law
unchanged, so every column sum is at most $D_X$. Next, $\Upsilon_Y$ and
$U_Y=r^2\sum_j2^{\varpi+|N_Y(j)|}(n-\frac12|N_Y(j)|)$ depend only on the multiset of
the column degrees, and $\Upsilon_X$, $U_X$ in the same way on the row degrees.
Finally $U_Y/r^2$ is a constant, so $(e^{U_Y}-1)/r^2$ is a power series in $r$ with
nonnegative coefficients; $e^{\rm pin}_d$, $\gamma_{\rm pin}$ and
$r^2n^22^{\Delta_{\rm pin}}$ are polynomials in $r$ with nonnegative coefficients; and
$\Upsilon_Y(1-e^{\rm pin}_\varpi)^{-n}$ and $\Lambda^{\rm pin}_0$ are built from them
by products, exponentials and the quotients $1/(1-\varphi)$ and
$(1+\varphi)/(1-\varphi)$. So \cref{lem:majorant}(i),(ii) give the admissibility.
\end{proof}

\paragraph{Star matrices.}
The star part of a pinned row is of first order, and the extraction does not touch
it. Under a product of coordinate laws the leaves supported on
disjoint exact supports (a leaf has \emph{exact support} $S$ if
it depends exactly on the coordinates in $S$ and has mean zero in
each of them) are orthogonal, and the complete star part factors
through one coefficient matrix; the
leaves are centred in their own laws, and their means, bounded by
$\beta_\chi$ in \cref{prop:conditional-star}(b), contribute the separate term
$T_{\rm mean}$ there.

\begin{proposition}[Star matrix under pinned product laws]\label{prop:conditional-star}
Let $p\equiv3\pmod4$, let $\mathcal H=L^2_0(\Fp)$ for the uniform law, and let
$\Phat=C/r$ be the normalized Paley operator, so that $\Phat\one=0$, $\Phat^*=-\Phat$
and $\Phat^2=-I$ on $\mathcal H$ (\cref{lem:spec:gauss}). Let $X$ and $Y$ be finite
sets, let $\nu_X=\bigotimes_{i\in X}\nu_i$ and $\nu_Y=\bigotimes_{j\in Y}\nu'_j$ be
product laws on $\Fp^X$ and $\Fp^Y$ whose coordinate laws have densities at most
$\kappa$ relative to the uniform law and are supported in a set $\Omega\subseteq\Fp$,
and let $D\in\{0,\pm1\}^{X\times Y}$, $N_X(i)=\{j:D_{ij}\ne0\}$ and
$N_Y(j)=\{i:D_{ij}\ne0\}$. In the expansion of the kernel
$\prod_{i,j}(1+D_{ij}\chi(y_j-x_i))$ the \emph{star part} consists of the single pairs
$D_{ij}\chi(y_j-x_i)$, the \emph{row stars} $\prod_{j\in T}D_{ij}\chi(y_j-x_i)$ with
centre $i$ and leaves $T$, and the \emph{column stars}
$\prod_{i\in S}D_{ij}\chi(y_j-x_i)$ with centre $j$ and leaves $S$ ($|T|,|S|\ge2$);
let $T_\star\colon L^2_0(\nu_Y)\to L^2_0(\nu_X)$ be its operator relative to
$\nu_X\otimes\nu_Y$, compressed to the mean-zero functions on both sides. (For a
pinned row at $p\equiv3\pmod4$, $D_{ij}=\varsigma_{ij}$ for
$(i,j)\in E_\times$, since $\chi(\varsigma_{ij})=\varsigma_{ij}$, and
$D_{ij}=0$ off $E_\times$.)
For sets $T\subseteq Y$ and $S\subseteq X$ with at least two elements put
$A_{iT}=(-1)^{|T|}\prod_{j\in T}D_{ij}$ and $B_{jS}=\prod_{i\in S}D_{ij}$, and let
$M(\Phat)$ be the block operator matrix with rows indexed by $X$ and the sets $S$ and
columns indexed by $Y$ and the sets $T$, with the singleton block $D\otimes\Phat$, the
blocks $A\otimes I_{\mathcal H}$ (rows $i$, columns $T$) and $B^*\otimes I_{\mathcal H}$
(rows $S$, columns $j$), and zero between sets, as in \cref{cor:star-matrix}; let
$M(i)=\begin{psmallmatrix}iD&A\\B^*&0\end{psmallmatrix}$.
\begin{enumerate}[label=\textup{(\alph*)}]
\item For $\mathsf k>0$, $\norm{M(i)}\le\mathsf k$ if and only if
 \begin{equation}\label{eq:conditional-star-gate}
  \begin{pmatrix}\mathsf k^2I-AA^*&\mathsf kD\\ \mathsf kD^*&\mathsf k^2I-BB^*\end{pmatrix}
  \succeq0 .
 \end{equation}
\item Let $\beta_\chi$ bound $|\E_{\nu_i}\chi(t-x_i)|$ and $|\E_{\nu'_j}\chi(t-y_j)|$ for
 all $i$, $j$ and $t\in\Omega$. Let $\mathsf d_{\rm L}$ and $\mathsf d_{\rm R}$ be the
 square roots of the largest row sum times the largest column sum of the nonnegative
 $X\times Y$ matrices with the entries $(1+\beta_\chi)^{|N_X(i)|-1}-1$, respectively
 $(1+\beta_\chi)^{|N_Y(j)|-1}-1$, where $D_{ij}\ne0$, and zero elsewhere, and
 $\mathsf d_1=\mathsf d_{\rm L}+\mathsf d_{\rm R}$; with $m_{ii'}=|N_X(i)\cap N_X(i')|$ put
 \[
  \mathsf e_{\rm L}^2=\max_{i\in X}\sum_{i'\in X}\sum_{s=2}^{m_{ii'}}\binom{m_{ii'}}s
  \bigl((1+\beta_\chi)^{|N_X(i)|-s}-1\bigr)\bigl((1+\beta_\chi)^{|N_X(i')|-s}-1\bigr),
 \]
 let $\mathsf e_{\rm R}$ be defined likewise with $N_Y$, and put
 $\mathsf e=\max(\mathsf e_{\rm L},\mathsf e_{\rm R})$ and
 $\mathsf e_\star=\frac12\bigl(\mathsf d_1+\sqrt{\mathsf d_1^2+4\mathsf e^2}\bigr)$. Then
 $T_\star=\kappa r\,\mathcal F_XM(\Phat)\mathcal F_Y^*+T_{\rm mean}$, with contractions
 $\mathcal F_X\colon\mathcal H^X\oplus\mathcal H^{\mathcal S_X}\to L^2_0(\nu_X)$,
 where $\mathcal S_X=\{S\subseteq X:|S|\ge2\}$, and $\mathcal F_Y$ likewise
 with $Y$ in place of $X$, built from $\nu_X$, respectively $\nu_Y$, alone
 (defined in the proof) and $\norm{T_{\rm mean}}\le\kappa r\mathsf e_\star$. Hence
 $\norm{T_\star}\le\kappa r(\norm{M(i)}+\mathsf e_\star)$, and under (a)
 $\norm{T_\star}\le\kappa r(\mathsf k+\mathsf e_\star)$. No invariance of the coordinate
 laws is used.
\item \emph{Decay.} Let $\kappa_0\ge\kappa$ and $0<\omega_s\le1$ ($s\ge2$) satisfy
 $\kappa_0\omega_s^2\ge\min\{\kappa_0,1+\nu_\chi^s/r^2\}$, where $\nu_\chi$ bounds
 $|\operatorname{Cov}_\nu(\chi(t-x),\chi(t'-x))|$ for every coordinate law $\nu$ and
 all distinct $t,t'\in\Omega$ (if $\omega_s=1$ for every $s\ge2$ no bound on
 $\nu_\chi$ is needed; an individual $s$ with $\omega_s<1$ uses the bound at
 that $s$). Let
 $A_\omega$, $B_\omega$ be $A$, $B$ with every column of a set $T$, respectively $S$,
 multiplied by $\omega_{|T|}$, respectively $\omega_{|S|}$, and let $M_\omega$ be built
 from them. Then (b) holds with $\kappa_0$, $M_\omega$ and contractions
 $\mathcal F^\omega_X$, $\mathcal F^\omega_Y$ in place of $\kappa$, $M$, $\mathcal F_X$,
 $\mathcal F_Y$; so \eqref{eq:conditional-star-gate} for $A_\omega$, $B_\omega$ gives
 $\norm{T_\star}\le\kappa_0r(\mathsf k+\mathsf e_\star)$. If $\kappa$ and $\nu_\chi/r$
 are admissible functions of $r$ (\cref{lem:majorant}) and rational $\kappa_0$,
 $\omega_s$ satisfy these conditions at a threshold $p_0$, they satisfy them at every
 prime $p\ge p_0$.
\end{enumerate}
\end{proposition}

\begin{proof}
\emph{Centred feature maps.} (The maps $U_\nu$ and $W_S$ below
are the \emph{feature maps}; $\mathcal F_X,\mathcal F_Y$ are
assembled from them.) For a coordinate law $\nu$ on $\Fp$ with
density $\rho\le\kappa$ let
$\iota_\nu\colon L^2(\Fp)\to L^2(\nu)$ be the evaluation $f\mapsto f$; then
$\norm{\iota_\nu f}^2_\nu=\E\rho |f|^2\le\kappa\norm f^2$ and $\iota_\nu^*g=\rho g$. The
\emph{centred evaluation} $U_\nu=P_0\iota_\nu|_{\mathcal H}\colon\mathcal H\to L^2_0(\nu)$
has norm at most $\sqrt\kappa$; since $P_0\iota_\nu\one=0$,
$P_0\iota_\nu=U_\nu P_{\mathcal H}$, with $P_{\mathcal H}$ the projection of $L^2(\Fp)$
onto $\mathcal H$; and for $g\in L^2_0(\nu)$, $U_\nu^*g=\rho g$ is a mean-zero function
supported in $\supp\nu\subseteq\Omega$. Put $m_\nu(t)=\E_{x\sim\nu}\chi(t-x)$ and
$\phi_\nu(t,x)=\chi(t-x)-m_\nu(t)$, so that $\E_{x\sim\nu}\phi_\nu(t,x)=0$. For
$S\subseteq X$ with $|S|\ge2$ let $W_S\colon L^2(\Fp)\to L^2(\nu_X)$,
$(W_Sf)(x)=\E_tf(t)\prod_{i\in S}\phi_{\nu_i}(t,x_i)$, and define $W'_T$ for
$T\subseteq Y$ likewise with the laws $\nu'_j$. Write $U_i=U_{\nu_i}$ and
$U'_j=U_{\nu'_j}$, their outputs read as functions of $x_i$ and of $y_j$.

\emph{Step 1 (the contractions $\mathcal F_X,\mathcal F_Y$).} The product law gives the orthogonal decomposition
$L^2(\nu_X)=\bigoplus_{R\subseteq X}H_R$, where $H_R$ consists of the functions of
$x_R$ whose conditional mean in every coordinate of $R$ vanishes; $U_i$ maps into
$H_{\{i\}}$ and $W_S$ into $H_S$. By independence,
$\ip{W_Sf}{W_Sf'}=\E_{t,t'}\overline{f(t)}f'(t')\prod_{i\in S}\Gamma_i(t,t')$ with
$\Gamma_i(t,t')=\E_{x\sim\nu_i}\phi_{\nu_i}(t,x)\phi_{\nu_i}(t',x)$. Each $\Gamma_i$ is a
positive semidefinite kernel with $\Gamma_i(t,t)\le1$, and as an operator relative to
the uniform law it is $V_i^*V_i$ with
$V_if=\E_t\phi_{\nu_i}(t,\cdot)f(t)=P_0\iota_{\nu_i}Cf$, so $0\preceq\Gamma_i\preceq\kappa
r^2I$. For $i_1\in S$, the pointwise product of the $\Gamma_i$ with $i\ne i_1$ is a
positive semidefinite kernel with diagonal at most one (Schur product theorem), so by
\cref{lem:swapped-rectangle}(b) the kernel $\prod_{i\in S}\Gamma_i$ defines an operator
between $0$ and $\kappa r^2I$, and $\norm{W_S}\le\sqrt\kappa\,r$ on all of $L^2(\Fp)$.
As the ranges are orthogonal,
\[
 \mathcal F_X\bigl((f_i)_{i\in X},(g_S)_S\bigr)=\sum_{i\in X}\frac{U_if_i}{\sqrt\kappa}
 +\sum_S\frac{W_Sg_S}{\sqrt\kappa\,r}
\]
is a contraction from $\mathcal H^X\oplus\mathcal H^{\mathcal S_X}$ into $L^2_0(\nu_X)$,
$\mathcal S_X$ the set of the $S\subseteq X$ with $|S|\ge2$; $\mathcal F_Y$ is built in
the same way from the $U'_j$ and the $W'_T$.

\emph{Step 2 (exact supports).} Let $g\in L^2_0(\nu_Y)$, and let $\bar g_j$ be its
component in $H_{\{j\}}$. A kernel has \emph{exact support} $R$ in $y$ if, for
fixed $x$, it depends only on $y_R$ and has conditional mean zero in each
coordinate of $R$. Such a kernel sees only the component of $g$ in $H_R$. \emph{Single pairs:}
$\E_y\chi(y_j-x_i)g(y)=\E_b\chi(b-x_i)\rho'_j(b)\bar g_j(b)=(CU'^*_j\bar g_j)(x_i)$, with
$U'^*_j\bar g_j\in\mathcal H$; after the output is centred, the pair $(i,j)$ gives
$rD_{ij}U_i\Phat U'^*_j$. \emph{Row stars:} since $\chi(-1)=-1$,
$\prod_{j\in T}D_{ij}\chi(y_j-x_i)=A_{iT}\prod_{j\in T}\bigl(m_{\nu'_j}(x_i)
+\phi_{\nu'_j}(x_i,y_j)\bigr)$. In the expansion, the term with all factors $\phi$ has
exact support $T$ and gives, after centring,
$A_{iT}P_0\iota_{\nu_i}W'^*_T=A_{iT}U_iP_{\mathcal H}W'^*_T=A_{iT}U_i(W'_T|_{\mathcal H})^*$;
the term without a factor $\phi$ is a function of $x_i$ and vanishes on mean-zero $g$;
the others, with exact supports $\emptyset\ne T'\subsetneq T$, carry the
\emph{complementary means} $\prod_{j\in T\setminus T'}m_{\nu'_j}(x_i)$. \emph{Column
stars:} $\prod_{i\in S}D_{ij}\chi(y_j-x_i)=B_{jS}\prod_{i\in S}\bigl(m_{\nu_i}(y_j)
+\phi_{\nu_i}(y_j,x_i)\bigr)$; the term with all factors $\phi$ gives
$B_{jS}W_SU'^*_j$, the term without a factor $\phi$ has a constant image, which the
centring removes, and the others carry complementary means. So
$T_\star=T_{\rm pure}+T_{\rm mean}$, where $T_{\rm mean}$ collects the terms with
complementary means and
\[
 T_{\rm pure}=\sum_{i,j}rD_{ij}U_i\Phat U'^*_j+\sum_{i,T}A_{iT}U_i(W'_T|_{\mathcal H})^*
 +\sum_{j,S}B_{jS}W_SU'^*_j=\kappa r\,\mathcal F_XM(\Phat)\mathcal F_Y^*,
\]
as one checks block by block.

\emph{Step 3 (the pure part).} By \cref{lem:spec:gauss},
$\mathcal H=\mathcal H_+\oplus\mathcal H_-$ with $\Phat=\pm iI$ on $\mathcal H_\pm$, and
$M(\Phat)$ acts on the two parts as $M(i)\otimes I$ and as $\overline{M(i)}\otimes I$;
so $\norm{M(\Phat)}=\norm{M(i)}$ and $\norm{T_{\rm pure}}\le\kappa r\norm{M(i)}$. The
matrix is decomposed before the contractions are applied, so no commutation of $\Phat$
with the features is used.

\emph{Step 4 (the gate).} For $\mathsf k>0$, $\norm{M(i)}\le\mathsf k$ if and only if
the Hermitian dilation $\begin{psmallmatrix}\mathsf kI&M(i)\\M(i)^*&\mathsf kI
\end{psmallmatrix}$ is positive semidefinite. Order its index set as
$(X,\mathcal S_X,Y,\mathcal T_Y)$, $\mathcal T_Y$ the sets $T$, and take the Schur
complement of its diagonal block $\mathsf kI$ on $(\mathcal S_X,\mathcal T_Y)$: the
dilation is positive semidefinite if and only if, by
\cref{lem:block-completion}(a),
\[
 \begin{pmatrix}\mathsf kI-\mathsf k^{-1}AA^*&iD\\-iD^*&\mathsf kI-\mathsf k^{-1}BB^*
 \end{pmatrix}\succeq0 .
\]
Multiplying by $\mathsf k$ and conjugating with the unitary $\diag(I,-iI)$ gives
\eqref{eq:conditional-star-gate}. This proves (a).

\emph{Step 5 (complementary means).} Fix a row centre $i$ and an exact support
$T'\ne\emptyset$. The terms of the row stars at $i$ with this support add up to
$(-1)^{|T'|}\prod_{j\in T'}D_{ij}\,c(x_i)\prod_{j\in T'}\phi_{\nu'_j}(x_i,y_j)$, where
$c(t)=\prod_{l\in N_X(i)\setminus T'}(1-D_{il}m_{\nu'_l}(t))-1$ satisfies
$|c|\le(1+\beta_\chi)^{|N_X(i)|-|T'|}-1$ on $\Omega$. Their operator maps
$H_{T'}(\nu_Y)$ into $H_{\{i\}}(\nu_X)$: for $|T'|\ge2$ it is
$g\mapsto\pm P_0[c\cdot\iota_{\nu_i}W'^*_{T'}g]$, of norm at most $\kappa r\norm c_\infty$
by Step 1, and for $T'=\{j\}$ it is $g\mapsto\pm P_0[c\cdot\iota_{\nu_i}C^*U'^*_j\bar
g_j]$, with the same bound, since
$\E_{y_j}\phi_{\nu'_j}(t,y_j)\bar g_j(y_j)=(C^*U'^*_j\bar g_j)(t)$. Likewise the column
stars at $j$ with exact support $S'$ give an operator from $H_{\{j\}}(\nu_Y)$ into
$H_{S'}(\nu_X)$, namely $g\mapsto W_{S'}\iota_{\nu'_j}^*(b\,\bar g_j)$ with
$|b|\le(1+\beta_\chi)^{|N_Y(j)|-|S'|}-1$ on $\Omega$, and for $S'=\{i\}$ the same with
$P_0\iota_{\nu_i}C$ in place of $W_{S'}$; its norm is at most $\kappa r\norm b_\infty$.
Write $L^2_0(\nu_X)=\mathcal V_X\oplus\mathcal W_X$ with
$\mathcal V_X=\bigoplus_iH_{\{i\}}$ and $\mathcal W_X=\bigoplus_{|S|\ge2}H_S$, and likewise
for $Y$. The norm of an operator matrix between orthogonal sums is at most the norm of
the matrix of the norms of its entries. The block $\mathcal V_X\leftarrow\mathcal V_Y$
of $T_{\rm mean}$ has entries $H_{\{i\}}\leftarrow H_{\{j\}}$ of norm at most $\kappa r$
times the sum of the entries of the two matrices of (b), so its norm is at most
$\kappa r(\mathsf d_{\rm L}+\mathsf d_{\rm R})$ by Schur's test. The block
$\mathcal V_X\leftarrow\mathcal W_Y$ has entries $H_{\{i\}}\leftarrow H_{T'}$ of norm at
most $\kappa re_{iT'}$, with $e_{iT'}=(1+\beta_\chi)^{|N_X(i)|-|T'|}-1$ for
$T'\subseteq N_X(i)$ and zero otherwise; the squared norm of the matrix $(e_{iT'})$ is
the norm of its Gram matrix $(\sum_{T'}e_{iT'}e_{i'T'})_{i,i'}$, at most its largest row
sum $\mathsf e_{\rm L}^2$. Likewise the block $\mathcal W_X\leftarrow\mathcal V_Y$ has
norm at most $\kappa r\mathsf e_{\rm R}$, and the block $\mathcal W_X\leftarrow\mathcal W_Y$
vanishes, since every star has a singleton centre. So
$\norm{T_{\rm mean}}\le\kappa r\norm{\begin{psmallmatrix}\mathsf d_1&\mathsf e\\
\mathsf e&0\end{psmallmatrix}}=\kappa r\mathsf e_\star$, which with Step 3 proves (b).

\emph{Step 6 (decay).} The Gram kernel of $W_S\one_\Omega$ is
$\one_\Omega(t)\one_\Omega(t')\prod_{i\in S}\Gamma_i(t,t')$, with diagonal at most one and
off-diagonal entries of modulus at most $\nu_\chi^{|S|}$; its absolute row sums relative
to the uniform law are at most $r^2+\nu_\chi^{|S|}$, and Schur's test gives
$\norm{W_S\one_\Omega}^2\le r^2+\nu_\chi^{|S|}$. Since $U'^*_j$ has mean-zero range
supported in $\Omega$, $W_SU'^*_j=W_S\one_\Omega P_{\mathcal H}U'^*_j$, the two
projections being applied in this order; and since $\iota_{\nu_i}$ sees only $\Omega$,
$U_iP_{\mathcal H}W'^*_T=U_iP_{\mathcal H}\one_\Omega W'^*_T=U_i(W'_T\one_\Omega
P_{\mathcal H})^*$. By the choice of $\omega_s$,
$\norm{W_S\one_\Omega P_{\mathcal H}}\le\sqrt{\kappa_0}\,r\omega_{|S|}$: if
$\omega_{|S|}<1$, then $\kappa_0\omega_{|S|}^2\ge1+\nu_\chi^{|S|}/r^2$, and otherwise
Step 1 applies. So the maps of Step 1 with $\sqrt{\kappa_0}$ in place of $\sqrt\kappa$
and $W_S\one_\Omega P_{\mathcal H}/(\sqrt{\kappa_0}\,r\omega_{|S|})$ in place of
$W_S/(\sqrt\kappa\,r)$, and likewise for $Y$, are contractions $\mathcal F^\omega_X$,
$\mathcal F^\omega_Y$, and $T_{\rm pure}=\kappa_0r\,\mathcal F^\omega_XM_\omega(\Phat)
\mathcal F^{\omega*}_Y$, while $\norm{T_{\rm mean}}\le\kappa r\mathsf e_\star\le\kappa_0r
\mathsf e_\star$. At a prime $p\ge p_0$, $\kappa(p)\le\kappa(p_0)\le\kappa_0$, and
$1+\nu_\chi^s/r^2=1+(\nu_\chi/r)^sr^{s-2}$ ($s\ge2$) is at most its value at $p_0$, by
admissibility; the gate does not depend on $p$.
\end{proof}

For an invariant law the same matrix appears in \cref{cor:star-matrix}; under a
product law the leaves are centred in their own laws and their means are paid by
$\mathsf e_\star$. In the conditional rows of the upper range
(\cref{subsec:push:layer72product,subsec:advance:product-rows}) the coordinate laws
are the allowed-set laws, $\Omega$ is the complement of the middle labels,
$\kappa=2^\varpi/(1-e^-_\varpi(z))$, $\beta_\chi$ is the one-character mean bound and $\nu_\chi$
the covariance bound (display (PR3) in
\cref{subsec:advance:product-rows}); in the bands the bound
$\min\{C_\star,
\kappa_0r(\mathsf k+\mathsf e_\star)\}$ takes the place of the scalar majorant $C_\star$
of \eqref{eq:row-anova-stars} in the correlation bound \eqref{eq:row-anova-actual}.

\paragraph{Joint fold and matching rows.}
The completion of an upper coordinate bounds its fold row and its matching row
jointly, relative to one chart (\cref{lem:five-type}(f)). The two rows have the same
allowed sets, so their star parts factor through the same contractions, and the joint
bound is one more matrix test.

\begin{lemma}[Joint fold and matching rows]\label{lem:joint-rows}
Let $p\equiv3\pmod4$, and in the setting of this subsection let $\kappa$ bound the
densities of the allowed-set laws $\lambda_i$ relative to the uniform law and
$I_0\le\lambda(\mathcal I_{\rm pin})$. Let two sets of signed cross pairs be given, the
\emph{fold} row $\rm F$ and the \emph{matching} row $\rm E$, with signed incidence
matrices $D_{\rm F},D_{\rm E}\in\{0,\pm1\}^{n\times n}$ ($D_\bullet(i,j)=\varsigma_{ij}$ on
the pairs of the row and zero elsewhere) and $D_{\rm E}^{\mathsf T}=-D_{\rm E}$. For
$\bullet\in\{{\rm F},{\rm E}\}$ let $K_\bullet$ be the raw kernel relative to
$\mu\otimes\mu$ (the indicator that $x([n])\cap y([n])=\emptyset$, times the product over
the pairs of $1+\chi_{ij}$), $Z_\bullet=\ip\one{K_\bullet\one}_\mu$ and
$\pi_\bullet=K_\bullet\,\mu\otimes\mu/Z_\bullet$. Let $\nu$ be the marginal of
$\pi_{\rm E}$, and suppose $d\nu/d\mu\ge e^{-\ell_{\rm E}}$. Relative to $\nu\otimes\nu$
let $T_{\rm F}$ and $T_{\rm E}$ be the densities of $\pi_{\rm F}$ and $\pi_{\rm E}$, and
$u_{\rm L}$, $u_{\rm R}$ the densities of the two marginals of
$\pi_{\rm F}$ relative to $\nu$, and put
$C_{\rm F}=T_{\rm F}-u_{\rm L}\otimes u_{\rm R}$ and $C_{\rm E}=T_{\rm E}-\Jproj$, as
operators on $L^2(\nu)$. With $\kappa_0$ and $\omega_s$ as in
\cref{prop:conditional-star}(c) for the laws $\lambda_i$ and
$\Omega=\Fp\setminus\mathcal M$, put, for $D=D_{\rm F},D_{\rm E}$ and all
$S\subseteq[n]$ with $|S|\ge2$,
\[
 A_D(i,S)=(-1)^{|S|}\omega_{|S|}\prod_{j\in S}D(i,j),\qquad
 B_D(i,S)=\omega_{|S|}\prod_{j\in S}D(j,i),
\]
and
\[
 \mathcal Z=\begin{pmatrix}D_{\rm E}&D_{\rm F}\\-D_{\rm F}^{\mathsf T}&D_{\rm E}\end{pmatrix},
 \qquad
 \mathcal A=\begin{pmatrix}A_{D_{\rm E}}&A_{D_{\rm F}}\\B_{D_{\rm F}}&A_{D_{\rm E}}\end{pmatrix},
 \qquad\mathcal G=\mathcal A\mathcal A^*.
\]
Then $C_{\rm F}$ and $C_{\rm E}$ annihilate the constants on both sides,
$C_{\rm E}=C_{\rm E}^*$, and $A_{D_{\rm E}}=B_{D_{\rm E}}$.
\begin{enumerate}[label=\textup{(\alph*)}]
\item Let $\mathsf k>0$ satisfy
 \begin{equation}\label{eq:joint-gate}
  \mathsf k^2I+i\mathsf k\mathcal Z-\mathcal G\succeq0,
 \end{equation}
 a Hermitian test of size $2n$, rational of size $4n$ in real form
 (\cref{rem:realified-gates}). For each row $\bullet$ let $e_\bullet$ bound the norm of
 the complementary-mean operator $T^\bullet_{\rm mean}$ of
 \cref{prop:conditional-star}(b) for $D_\bullet$ on $L^2_0(\lambda)$, where
 $\lambda$ is the product of the allowed-set laws, divided by $\kappa_0r$
 (for instance $\mathsf e_\star$ of \cref{prop:conditional-star}(b) for
 $D_\bullet$); let $S^\star_\bullet\ge\norm{P_0K^{\rm st}_\bullet P_0}$ for the star part
 $K^{\rm st}_\bullet$ of $K_\bullet$, $N_\bullet\ge\norm{T^\bullet_{\rm ns}}
 +\norm{T^\bullet_{\rm coll}}$ for its nonstar and collision parts,
 $\eps_\bullet\ge|Z_\bullet-1|$ with $\eps_\bullet<1$, and
 $m_\bullet\ge\norm{P_0K_\bullet\one},\norm{P_0K_\bullet^*\one}$, all on $L^2(\mu)$. Then
 \[
  \begin{pmatrix}C_{\rm E}&C_{\rm F}\\C_{\rm F}^*&C_{\rm E}\end{pmatrix}
  \succeq-\mathfrak k_pI\quad\text{on }L^2(\nu)\oplus L^2(\nu),\qquad
  \mathfrak k_p=e^{\ell_{\rm E}}\Bigl[\frac{\kappa_0r(\mathsf k+e_{\rm F}+e_{\rm E})}{I_0}
  +\mathfrak c_{\rm F}+\mathfrak c_{\rm E}\Bigr],
 \]
 where $\mathfrak c_\bullet=(\eps_\bullet S^\star_\bullet+N_\bullet)/(1-\eps_\bullet)
 +m_\bullet^2/(1-\eps_\bullet)^2$.
\item Let $\Blk$ be a finite set of primes and, for $p\in\Blk$, let $C_{{\rm F},p}$ and
 $C_{{\rm E},p}=C_{{\rm E},p}^*$ be operators on $L^2(\nu_p)$ that annihilate the
 constants on both sides, for instance those of (a) at $p$. Put
 $C_\Blk=\bigotimes_p(\Jproj_p+C_{{\rm F},p})-\Jproj_\Blk$ and
 $E_\Blk=\bigotimes_p(\Jproj_p+C_{{\rm E},p})-\Jproj_\Blk$ on $L^2(\bigotimes_p\nu_p)$.
 They act on each exact prime-support sector
 $\bigotimes_{p\in S}L^2_0(\nu_p)\otimes\one$ ($\emptyset\ne S\subseteq\Blk$) as
 $\bigotimes_{p\in S}C_{{\rm F},p}$ and $\bigotimes_{p\in S}C_{{\rm E},p}$, so the joint
 block of $(E_\Blk,C_\Blk)$ is the direct sum of the joint blocks of the sectors. If
 $\norm{C_{{\rm F},p}}\le c<1$ and $\norm{C_{{\rm E},p}}\le h<1$ for every $p$, the
 joint block of every sector with at least two primes is $\succeq-(c^2+h^2)I$; so the
 joint block of $(E_\Blk,C_\Blk)$ is $\succeq-\max(c^2+h^2,\max_p\mathfrak k'_p)I$
 whenever the joint block at each prime $p$ is $\succeq-\mathfrak k'_pI$.
\item If moreover $C_{{\rm F},p}=T_{{\rm F},p}-u_{{\rm L},p}\otimes u_{{\rm R},p}$ for a
 kernel $T_{{\rm F},p}$ of mass one with marginal densities $u_{{\rm L},p}$,
 $u_{{\rm R},p}$ relative to $\nu_p$, and $U_{\rm L}=\bigotimes_pu_{{\rm L},p}$,
 $U_{\rm R}=\bigotimes_pu_{{\rm R},p}$ satisfy
 $\norm{U_{\rm L}-\one},\norm{U_{\rm R}-\one}\le\upsilon$, then
 $T_{{\rm F},\Blk}=\bigotimes_pT_{{\rm F},p}$ satisfies
 \[
  \norm{T_{{\rm F},\Blk}-U_{\rm L}\otimes U_{\rm R}-C_\Blk}\le(2\upsilon+\upsilon^2)
  \Bigl[\prod_{p\in\Blk}\bigl(1+\norm{C_{{\rm F},p}}\bigr)-1\Bigr],
 \]
 which is at most $(2\upsilon+\upsilon^2)(e^{Nc}-1)$ if $\Blk$ has at most $N$
 primes, for an integer $N\ge|\Blk|$, and $\norm{C_{{\rm F},p}}\le c$ for
 every $p\in\Blk$.
\end{enumerate}
\end{lemma}

\begin{proof}
\emph{Preliminaries.} At $p\equiv3\pmod4$ the factor of a pair is
$1+\varsigma_{ij}\chi(y_j-x_i)$. The factor of $(i,j)$ in $K_{\rm E}(y,x)$ is
$1+\varsigma_{ij}\chi(x_j-y_i)=1+D_{\rm E}(j,i)\chi(y_i-x_j)$, the factor of $(j,i)$ in
$K_{\rm E}(x,y)$; so $K_{\rm E}$ is symmetric, $\pi_{\rm E}$ has two equal marginals
$\nu$, and $T_{\rm E}$ is a symmetric kernel with marginal densities one; hence
$C_{\rm E}=C_{\rm E}^*$ annihilates the constants. As $T_{\rm F}\one=u_{\rm L}$ and
$\E_\nu u_{\rm R}=1$, $C_{\rm F}\one=0$, and likewise $C_{\rm F}^*\one=0$. Finally
$B_{D_{\rm E}}(i,S)=\omega_{|S|}\prod_{j\in S}(-D_{\rm E}(i,j))=A_{D_{\rm E}}(i,S)$.

(a) \emph{The actual chart.} Let $g=d\nu/d\mu=K_{\rm E}\one/Z_{\rm E}$,
$g_{\rm L}=K_{\rm F}\one/Z_{\rm F}$ and $g_{\rm R}=K_{\rm F}^*\one/Z_{\rm F}$. Then
$T_{\rm F}(x,y)=K_{\rm F}(x,y)/(Z_{\rm F}g(x)g(y))$, $u_{\rm L}=g_{\rm L}/g$ and
$u_{\rm R}=g_{\rm R}/g$, so for $f,h\in L^2(\nu)$
\[
 \ip f{C_{\rm F}h}_\nu=\E_{\pi_{\rm F}}[\overline{f(x)}h(y)]-\E_\nu[\overline f u_{\rm L}]\,\E_\nu[u_{\rm R}h]
 =\ip f{(K_{\rm F}/Z_{\rm F}-g_{\rm L}\otimes g_{\rm R})h}_\mu,
\]
and likewise $\ip f{C_{\rm E}h}_\nu=\ip f{(K_{\rm E}/Z_{\rm E}-g\otimes g)h}_\mu$. So the
quadratic form of the joint block at $(f,h)$ is that of the self-adjoint operator
$\mathbb K$ on $L^2(\mu)^2$ with the diagonal entries $K_{\rm E}/Z_{\rm E}-g\otimes g$
and the off-diagonal entries $K_{\rm F}/Z_{\rm F}-g_{\rm L}\otimes g_{\rm R}$ and its
adjoint, and $\norm f^2_\mu=\E_\nu[|f|^2/g]\le e^{\ell_{\rm E}}\norm f^2_\nu$. Hence
$\mathbb K\succeq-cI$ with $c\ge0$ gives the joint block $\succeq-ce^{\ell_{\rm E}}I$. A
block operator $\begin{psmallmatrix}X&Y\\Y^*&X\end{psmallmatrix}$ has norm at most
$\norm X+\norm Y$, which bounds every error term below.

\emph{Centring and normalization.} For a kernel $K$ with $Z=\ip\one{K\one}$ and
marginal densities $g_1=K\one/Z$, $g_2=K^*\one/Z$ the exact identity
$K/Z-g_1\otimes g_2=P_0KP_0/Z-(g_1-\one)\otimes(g_2-\one)$ holds, and for $K=K_\bullet$
the rank-one term has norm at most $m_\bullet^2/(1-\eps_\bullet)^2$. On
$\mathcal I_{\rm pin}\times\mathcal I_{\rm pin}$,
$K_\bullet=\one\otimes\one+K^{\rm st}_\bullet+T^\bullet_{\rm ns}-T^\bullet_{\rm coll}$,
the expansion of \cref{lem:spec:expansion} with the collision part of
\cref{prop:pinned-collision}; so
$P_0K_\bullet P_0/Z_\bullet=P_0K^{\rm st}_\bullet P_0+(Z_\bullet^{-1}-1)P_0K^{\rm st}_\bullet
P_0+Z_\bullet^{-1}P_0(T^\bullet_{\rm ns}-T^\bullet_{\rm coll})P_0$, and the last two terms
have norm at most $(\eps_\bullet S^\star_\bullet+N_\bullet)/(1-\eps_\bullet)$. Hence
$\mathbb K\succeq\mathbb K^{\rm st}-(\mathfrak c_{\rm F}+\mathfrak c_{\rm E})I$, where
$\mathbb K^{\rm st}$ is the joint block of the operators $P_0K^{\rm st}_\bullet P_0$.

\emph{Restriction.} Both rows have the law $\mu=\lambda(\cdot\mid\mathcal I_{\rm pin})$,
and $P=\lambda(\mathcal I_{\rm pin})\ge I_0$. By \cref{lem:restriction}(a), with the
isometry $\iota f=\one_{\mathcal I_{\rm pin}}f/\sqrt P$, which preserves centring,
$\ip F{\mathbb K^{\rm st}F}=\ip{\iota F}{\mathbb T\iota F}/P$ for $F\in L^2_0(\mu)^2$,
where $\mathbb T$ is the joint block, on $L^2_0(\lambda)^2$, of the operators
$T^\bullet_\star$ of \cref{prop:conditional-star} for $\nu_X=\nu_Y=\lambda$ and
$D=D_\bullet$; and $\mathbb K^{\rm st}$ vanishes on the constants. So
$\mathbb T\succeq-c'I$ gives $\mathbb K^{\rm st}\succeq-(c'/I_0)I$.

\emph{Pure stars.} By \cref{prop:conditional-star}(b),(c),
$T^\bullet_\star=\kappa_0r\,\mathcal FM_\bullet(\Phat)\mathcal F^*+T^\bullet_{\rm mean}$,
with one contraction $\mathcal F=\mathcal F^\omega_X=\mathcal F^\omega_Y$ for both rows and
both laws (it depends only on $\lambda$; take all sets $S\subseteq[n]$, $|S|\ge2$),
with $M_\bullet$ built from $D_\bullet$, $A_{D_\bullet}$ and $B_{D_\bullet}$, and with
$\norm{T^\bullet_{\rm mean}}\le\kappa_0re_\bullet$. So
$\mathbb T=\kappa_0r(\mathcal F\oplus\mathcal F)\mathbb M(\mathcal F\oplus\mathcal F)^*
+\mathbb T_{\rm mean}$ with $\norm{\mathbb T_{\rm mean}}\le\kappa_0r(e_{\rm E}+e_{\rm F})$,
where $\mathbb M$ has the diagonal entries $M_{\rm E}(\Phat)$ and the off-diagonal entries
$M_{\rm F}(\Phat)$ and $M_{\rm F}(\Phat)^*$. As $\Phat^*=-\Phat$, $M_{\rm F}(\Phat)^*$ has the
singleton block $-D_{\rm F}^{\mathsf T}\otimes\Phat$, the block $B_{D_{\rm F}}\otimes I$
from the sets to the singletons and $A_{D_{\rm F}}^*\otimes I$ from the singletons to the
sets; and $D_{\rm E}^{\mathsf T}=-D_{\rm E}$, $A_{D_{\rm E}}=B_{D_{\rm E}}$ give
$M_{\rm E}(\Phat)^*=M_{\rm E}(\Phat)$. Grouping the singletons of the two copies and the
sets of the two copies,
\[
 \mathbb M=\begin{pmatrix}\mathcal Z\otimes\Phat&\mathcal A\otimes I\\
 \mathcal A^*\otimes I&0\end{pmatrix}.
\]
On $\mathcal H_\pm$ it acts as $\begin{psmallmatrix}\pm i\mathcal Z&\mathcal A\\
\mathcal A^*&0\end{psmallmatrix}\otimes I$; these are complex conjugate Hermitian
matrices, $\mathcal Z$ being real and antisymmetric. By
\cref{lem:block-completion}(a), eliminating the block $\mathsf kI$, the one with $+i$ is $\succeq-\mathsf kI$ if and only if
$\mathsf kI+i\mathcal Z-\mathsf k^{-1}\mathcal G\succeq0$, which is \eqref{eq:joint-gate}.
So $\mathbb M\succeq-\mathsf kI$ and
$\mathbb T\succeq-\kappa_0r(\mathsf k+e_{\rm F}+e_{\rm E})I$, and the three steps give
$\mathfrak k_p$.

(b) As $C_{{\rm F},p}$ annihilates the constants on both sides,
$\bigotimes_p(\Jproj_p+C_{{\rm F},p})=\sum_{S\subseteq\Blk}\bigotimes_{p\in S}C_{{\rm F},p}
\otimes\bigotimes_{p\notin S}\Jproj_p$, the term of $S$ acting on the sector of $S$ and
vanishing on the others; the same holds for $E_\Blk$. On a sector with $|S|\ge2$ the
joint block is
$\succeq-(\norm{\bigotimes_{p\in S}C_{{\rm E},p}}+\norm{\bigotimes_{p\in S}C_{{\rm F},p}})I
\succeq-(h^{|S|}+c^{|S|})I\succeq-(h^2+c^2)I$; on a singleton sector it is the joint
block at the prime.

(c) Expand $T_{{\rm F},\Blk}=\bigotimes_p(u_{{\rm L},p}\otimes u_{{\rm R},p}+C_{{\rm F},p})$
over the sets $S\subseteq\Blk$ of the primes at which the factor $C_{{\rm F},p}$ is taken.
The term of $S=\emptyset$ is $U_{\rm L}\otimes U_{\rm R}$, and $C_\Blk$ is the sum of the
other terms with the marginal rank
$\bigotimes_{p\notin S}u_{{\rm L},p}\otimes u_{{\rm R},p}=a\otimes b$ of the inactive
primes replaced by $\Jproj=\one\otimes\one$. Since $\E u_{{\rm L},p}=1$,
$\norm{\bigotimes_{p\in A}u_{{\rm L},p}-\one}^2=\prod_{p\in A}\norm{u_{{\rm L},p}}^2-1$ is
nondecreasing in $A$, so $\norm{a-\one}\le\upsilon$, and likewise $\norm{b-\one}\le\upsilon$;
hence $a\otimes b-\one\otimes\one=(a-\one)\otimes b+\one\otimes(b-\one)$ has norm at most
$\upsilon(1+\upsilon)+\upsilon$. Summing $\prod_{p\in S}\norm{C_{{\rm F},p}}$ times this over
$S\ne\emptyset$ gives the bound, and $1+x\le e^x$ the last one.
\end{proof}

\begin{corollary}[The joint fluctuation in a positive row]\label{cor:joint-row-split}
In \cref{lem:joint-rows}(b) let $h\ge0$ and $\mathsf k_\Blk\ge h$ satisfy
$E_\Blk\succeq-hI$ and $\begin{psmallmatrix}E_\Blk&C_\Blk\\C_\Blk^*&E_\Blk\end{psmallmatrix}
\succeq-\mathsf k_\Blk I$ (a bound $\mathsf k$ of the joint block with $\mathsf k<h$ is
replaced by $h$). If $0\le m\le H$ and $a_{\rm S}\ge Ha_{\rm D}\ge0$, then
\[
 \begin{pmatrix}a_{\rm S}E_\Blk&ma_{\rm D}C_\Blk\\ma_{\rm D}C_\Blk^*&a_{\rm S}E_\Blk
 \end{pmatrix}\succeq-\bigl(ha_{\rm S}+Ha_{\rm D}(\mathsf k_\Blk-h)\bigr)I .
\]
With $T_{{\rm F},\Blk}-U_{\rm L}\otimes U_{\rm R}$ in place of $C_\Blk$ the bound decreases
by at most $Ha_{\rm D}(2\upsilon+\upsilon^2)[\prod_p(1+\norm{C_{{\rm F},p}})-1]$, by
\cref{lem:joint-rows}(c); the $S=\emptyset$ term
$U_{\rm L}\otimes U_{\rm R}$ of \cref{lem:joint-rows}(c) is not counted in
this fluctuation bound; it is instead absorbed by the mean allowance term
in \cref{lem:five-type}(f),
where $m$ is the fold mass $m_b$ and $a_{\rm S}$, $a_{\rm D}$ are the five-type weights.
\end{corollary}

\begin{proof}
The matrix is the sum of $ma_{\rm D}$ copies of the joint block and of
$a_{\rm S}-ma_{\rm D}$ copies of $\diag(E_\Blk,E_\Blk)$, both coefficients being
nonnegative as $m\le H$ and $a_{\rm S}\ge Ha_{\rm D}$. Its lower bound is
$-(ma_{\rm D}\mathsf k_\Blk+(a_{\rm S}-ma_{\rm D})h)=-(ha_{\rm S}+ma_{\rm D}(\mathsf
k_\Blk-h))$, and $m\le H$ is used only after $\mathsf k_\Blk-h\ge0$. The last assertion
is \cref{lem:joint-rows}(c), multiplied by $ma_{\rm D}\le Ha_{\rm D}$ in the off-diagonal
blocks.
\end{proof}

When the pins of a row make some allowed sets disjoint, the coordinates whose allowed
sets contain a given label lie in one member of the family $\mathcal S$ of
the lemma, and the density bound $\kappa_0$ then enters through the
principal submatrices on the members of $\mathcal S$.

\begin{lemma}[Positive potentials]\label{lem:joint-potential}
In \cref{lem:joint-rows}(a) let $\rho_i$ be the density of $\lambda_i$ relative to the
uniform law, and let $\mathcal S$ be a family of sets of coordinates $(c,i)$ of
$\mathcal Z$ ($c\in\{1,2\}$ the copy, $i\in[n]$) such that for every $t\in\Fp$ the
coordinates with $\rho_i(t)>0$ lie in one member of $\mathcal S$. Let $T$ be a real
symmetric $2n\times2n$ matrix with $T\pm i\mathcal Z\succeq0$, for instance
$T=|\mathcal Z|=(-\mathcal Z^2)^{1/2}$. If $\mathsf k>0$ and
\begin{equation}\label{eq:joint-potential-gate}
 \mathsf k^2I_S-\mathsf k\kappa_0T_S-\kappa_0^2\mathcal G_S\succeq0\qquad(S\in\mathcal S),
\end{equation}
where $X_S$ denotes the principal submatrix of $X$ on $S$, then the conclusion of
\cref{lem:joint-rows}(a) holds without \eqref{eq:joint-gate} and with $r\mathsf k/I_0$
in place of $\kappa_0r\mathsf k/I_0$ in $\mathfrak k_p$; all other terms are kept.
For rational $T$ and rational gate parameters, both
$T\pm i\mathcal Z\succeq0$ and \eqref{eq:joint-potential-gate} are rational tests in real
form (\cref{rem:realified-gates}).
\end{lemma}

\begin{proof}
In the proof of \cref{lem:joint-rows}(a) only the bound of the pure part
$\mathbb P=\kappa_0r(\mathcal F\oplus\mathcal F)\mathbb M(\mathcal F\oplus\mathcal F)^*$
changes. Let $\Psi_i\colon L^2(\lambda_i)\to L^2(\Fp)$, $g\mapsto\sqrt{\rho_i}\,g$, an
isometry onto the functions supported in the allowed set of $i$; then
$\Psi_i\iota_{\lambda_i}$ is the multiplication by $\sqrt{\rho_i}$, and
$U_{\lambda_i}=\Psi_i^*\Pi_i\sqrt{\rho_i}\,|_{\mathcal H}$, where $\Pi_i=\Psi_iP_0\Psi_i^*$ is
an orthogonal projection. Let $D_\rho$ be the multiplication by $\sqrt{\rho_i}$ on the
copy of $L^2(\Fp)$ of the coordinate $(c,i)$, extend $\Phat$ by zero to the constants,
and let $\mathcal W$ be the higher part of $\mathcal F\oplus\mathcal F$, a contraction.
Since the singleton part of $\mathcal F$ is $U_{\lambda_i}/\sqrt{\kappa_0}$, $\mathbb P$
is the compression, by $\Pi$ on the singleton coordinates (followed by $\Psi^*$), of
\[
 \mathbb S=\begin{pmatrix}rD_\rho(\mathcal Z\otimes\Phat)D_\rho&
 \sqrt{\kappa_0}\,rD_\rho(\mathcal A\otimes I)\mathcal W^*\\
 \sqrt{\kappa_0}\,r\mathcal W(\mathcal A^*\otimes I)D_\rho&0\end{pmatrix},
\]
and compressions keep lower bounds; so it suffices that $\mathbb S\succeq-r\mathsf kI$. By
\cref{lem:block-completion}(a) on the block $r\mathsf kI$ and
$\mathcal W^*\mathcal W\preceq I$,
this follows from $r\mathsf kI+rD_\rho(\mathcal Z\otimes\Phat)D_\rho-(\kappa_0r/\mathsf
k)D_\rho(\mathcal G\otimes I)D_\rho\succeq0$. On the eigenspaces of $\Phat$ for $i$, $-i$
and $0$ the operator $T\otimes I+\mathcal Z\otimes\Phat$ is $T+i\mathcal Z$, $T-i\mathcal Z$
and $T$, all positive semidefinite (the last by averaging the first two); so
$D_\rho(\mathcal Z\otimes\Phat)D_\rho\succeq-D_\rho(T\otimes I)D_\rho$, and it suffices that
$\mathsf k^2I-D_\rho((\mathsf kT+\kappa_0\mathcal G)\otimes I)D_\rho\succeq0$. This operator
is the multiplication by the matrix function $t\mapsto\mathsf k^2I-D(t)(\mathsf kT
+\kappa_0\mathcal G)D(t)$, $D(t)=\diag(\sqrt{\rho_i(t)})$. At each $t$, $D(t)$ vanishes
outside a set $S\in\mathcal S$ and $0\preceq D(t)\preceq\sqrt{\kappa_0}I$; as
$\mathsf kT+\kappa_0\mathcal G\succeq0$, the matrix $D(t)(\mathsf kT+\kappa_0\mathcal G)D(t)$
is at most $\kappa_0\lambda_{\max}((\mathsf kT+\kappa_0\mathcal G)_S)I\preceq\mathsf k^2I$,
by \eqref{eq:joint-potential-gate}. For $T=|\mathcal Z|$, $T\pm i\mathcal Z\succeq0$ because
$i\mathcal Z$ is Hermitian with $|i\mathcal Z|=|\mathcal Z|$.
\end{proof}

\begin{remark}[The pinned rows of the construction]\label{rem:pinned-instances}
The conditional rows of \cref{subsec:advance:product-rows} come in six geometries
(the other type L, F or E, with one or two layers), each for the two laws F and E.
In all twelve there are no edges inside a row, $n=z\in\{3,6,12\}$, and the pin bound
is $\varpi\in\{0,1,2,4\}$; in the two geometries with other type F the actual pin
numbers are mixed ($\{0,1\}$ and $\{0,2\}$) and are majorized by $\varpi$. The classes are those of
the compressed partition, grouped with $\mathfrak q=0$: the integer
$\mathfrak b=\mathrm{oc}(\mathcal C)$, at most $1280$, is used exactly, and for the
classes of case (A) of \cref{lem:advance:compressed-partition}(b) the better of two
covers is taken class by class (a class may carry several records,
\cref{def:extraction:records}). The E law of the E geometries has no cross pairs, so
its nonstar part vanishes; the other ten rows have $7164$ classes, counted with their
kernel constants $c_{\rm ker}(\mathcal C)$. In all twelve the multisets of row and
column degrees coincide. Hence $\Delta_{\rm pin}=\varpi+\max_i|N_X(i)|$ and
$D_X=D_Y$, and the evaluator, which computes
$\Upsilon_X(1-e^{\rm pin}_\varpi)^{-n}(e^{U_Y}-1)$, returns this common value. These
facts are read off the row geometry and the stored partitions
(\path{research/advance_2026_10_17/sparse/product_rows_all.py} and
\path{product_rows_all_parts.json}; the exact grouping is
\path{research/advance_2026_10_24/upper/conditional_fine_parts.json}, checked by
\path{check_conditional_fine_buckets.py} of the same directory into
\path{conditional_fine_parts_review.json}). Rows with internal edges, which would need
the merge charges of \cref{lem:matching:extractionk}, do not occur. The fold and
matching rows of one geometry have the same pins, which is the hypothesis of
\cref{lem:joint-rows}.
\end{remark}

\paragraph{The forms that the certificates evaluate.}
\emph{Nonstar bounds.} For the two root cuts of the upper range ($n=36$,
$\Delta^+=10$, $\Delta=20$) the partition is a refinement of the compressed
partition with witnesses of sizes $2$, $3$ and $4$, and the nonstar bound is
\cref{cor:extraction-leaf}(b) with the graded errors, the free partners of the
records (\cref{def:extraction:records}) and the constants of \eqref{eq:leaf-factors},
under the monotone-factor condition \eqref{eq:monotone-signature}. For the direct
Route~B criterion (\cref{lem:record:direct-star-inputs}), from $\Pthree$ to
$2^{52}$, the coefficients are evaluated at the nineteen points $2^{47},\dots,2^{52}$,
$167\cdot2^{40}$ and $j2^{42}$, $j\in\{40,42,43,44,46,48,50,52,56,60,80,96\}$, and
between consecutive ones bounded by \cref{lem:logconvex-coefficients}. The
Route~A criterion from $2^{52}$ on uses the tables $C_k(Q)$ at
$Q=2^{50},2^{52},\dots,2^{60}$, each prime using the table of the largest $Q\le p$;
these tables are computed with the coarse errors and the constants
$\max_i\sum_jw^{\rm m}_{ij}$, $\max_i\sum_jw^\varrho_{ij}$, which dominate the factors of
\eqref{eq:leaf-factors}.
The classes are stored as a histogram of signatures (the witness size,
$\mathrm{ex}^\star$, and the multiset of the pairs (kind, $s_v$), the kind being
$\mathsf f_{\rm m}$ on $U_x\cup U_z$ and $\mathsf f_\varrho$ on $C_x\cup C_z$) with integer
multiplicities, $30\,645\,955$ distinct signature rows in all ($17\,286\,284$ for LL20 and $13\,359\,671$
for LL22); every square root of a vertex factor is rounded upward to a multiple of
$2^{-48}$, and the integer sums over the histogram are exact. The nonroot caps
\ref{it:record:upper-F6} and the route-A twelve-label caps (\cref{sec:middle-reference})
use \cref{cor:extraction-compressed} with $\mathfrak q=0$ on the compressed partition,
each vertex summand of $\mathrm{oc}(\mathcal C)$ rounded upward at $2^{-20}$. The
conditional rows \ref{it:record:upper-F3} of the bands use
\cref{prop:pinned-extraction} with $\mathfrak q=0$ (\cref{rem:pinned-instances}),
\cref{prop:pinned-collision} and the star matrices of
\cref{prop:conditional-star}(a)--(c), and their joint fold and matching bounds
\cref{lem:joint-rows} (for the $\rm L$ rows with \cref{lem:joint-potential}); those of
the tail datum
(\cref{def:record:upper-bands}) are computed with $\mathfrak q=20$ and the scalar
majorant \eqref{eq:row-anova-stars}.
Soft kernels (\cref{cor:soft-kernels}) are evaluated with every
coefficient $D_e$ replaced by $|D_e|$ on its support
$\{e:D_e\ne0\}$ (its \emph{Boolean support}), or on a larger set
of pairs containing that support whose admissibility is checked
in \cref{app:certificates} (a \emph{certified Boolean
supergraph}).
\emph{Coincidence bounds.} For the root cuts the Route~B criterion takes
the minimum in \cref{lem:oct24-collision}(e) of $D_1+D_{\ge2}$ (with one order for
each pair $(i,j)$, one order $\pi_k$ for each $k=2,\dots,9$ chosen at each of the
points $2^{47},\dots,2^{52}$, and the tail (c) from $k_0=10$), $D_{\rm ord}$ and
$D_{\rm nat}$; a band evaluates this minimum at its start $X$, with the orders chosen
at the evaluation point of its band (\cref{cert:upper-root-caps}). The Route~A criterion takes
$\min(D_{\rm ord},D_{\rm nat})$. The nonroot caps and the route-A twelve-label caps use
$D_{\rm nat}$ of \eqref{eq:collision-natural}, the L20 cut of the middle root the
balanced bound of \cref{prop:l20-collisions} (\cref{cor:incident-weights}),
sharpened, for the rows in which the principal (root) law is the
uniform law, by the protected-subspace compression of
\cref{imp:local:collision-gap}, and the
product rows \cref{prop:pinned-collision}.

\paragraph{Consequences and what to read next.}
This section uses the labellings and the extension count with its coarse
and graded errors
(\cref{lem:extension}, \eqref{eq:graphs:et}, \eqref{eq:graphs:graded},
\ref{it:et:mono}--\ref{it:et:poly}); one-step error pairs, sequential counts, vertex
densities and marginal ratios
(\cref{def:graphs:error-pair,lem:graphs:sequential,lem:graphs:vertex}); spectral data,
their elementary symmetries and the cut kernel
(\cref{def:spec:data,lem:spec:elementary,lem:spec:cutkernel}); the positive Schur
order (\cref{lem:swapped-rectangle}(b)) and the restriction to an injectivity event
(\cref{lem:restriction}(a)); the character expansion, the Gauss operator and the real
form of the phase gates (\cref{lem:spec:expansion,lem:spec:gauss,rem:realified-gates}); the
four-cycle kernel for uniform probability (\cref{cor:dl:normalized}); and admissible
majorants (\cref{lem:majorant}). Three standard facts are used without statement: the
Schur product theorem, Schur's test for the norm of a kernel or a matrix, and
H\"older's inequality.
